\documentclass[12pt]{article}
\usepackage{titlesec}
\usepackage{xr}
\usepackage{array}
%\externaldocument{appendix,appendixtables,conclusion,data,figstables,intro_martin_v5,iv,literature,model_new2,panelresults,t1pa,t1pb,t1pc,t1pd,t2pa,t2pb,Table2,Table3,Table4,Table5,Table6,Table7,TableA1new,maininternetappendix}
\usepackage{endnotes}
\usepackage{indentfirst}
\usepackage{eurosym}
\usepackage[T1]{fontenc}
\usepackage[utf8]{inputenc}
%\DeclareUnicodeCharacter{00B4}{\'{}}
\usepackage{geometry}
\geometry{verbose,tmargin=1.1in,bmargin=1.1in,lmargin=.8in,rmargin=.8in}
\usepackage[active]{srcltx}
\usepackage{bbm}
\usepackage{amsthm}
\usepackage{amsmath}
\usepackage{setspace}
\usepackage{lscape}
\usepackage{multirow}
\usepackage{float}
%\usepackage[authoryear,longnamesfirst]{natbib}
\usepackage[authoryear]{natbib}
\setlength{\bibsep}{0.0pt}
\usepackage[inline]{enumitem}
\usepackage[T1]{fontenc}
\usepackage{graphicx}
\usepackage{tikz}
\usepackage{afterpage}
\usepackage{placeins}
\usepackage[labelfont=bf,labelsep=period,font=footnotesize,width=\textwidth,justification=justified,singlelinecheck=off]{caption}
\usepackage[font=footnotesize,labelformat=simple]{subcaption}
\renewcommand\thesubfigure{(\Alph{subfigure})}
\renewcommand\thesubtable{\thefigure, Panel \Alph{subtable}}
\DeclareCaptionLabelFormat{subtablelabel}{(\Alph{subtable})}
%\captionsetup[subtable]{labelformat=subtablelabel}
\newcommand{\RN}[1]{\textup{\uppercase\expandafter{\romannumeral#1}}}
\newcommand\inserttable[1]{\begin{center} \textbf{{\color{black}[Insert #1 Here.]}} \end{center}}
\usepackage{lmodern}
\usepackage{bbm}
\usepackage{relsize}
\global\long\def\sym#1{\ifmmode^{#1}\else\ensuremath{^{#1}}\fi}
\setstretch{1.5}
\makeatletter

\newtheorem{proposition}{Proposition}
\newtheorem{corollary}{Corollary}
\newtheorem{theorem}{Theorem}
\DeclareMathOperator*{\argmax}{arg\,max}
\usepackage{amsfonts}\usepackage[page]{appendix}
\usepackage{color}
%\usepackage[usenames,dvipsnames,svgnames,table]{xcolor}

\usepackage[hyphens]{url}
\usepackage{hyperref}
\hypersetup{
%	unicode=true,
	breaklinks=true,
	pdfborder={0 0 1},
	pdftitle={Innovation: The Bright Side of Common Ownership?},
	pdfauthor={Ant\'{o}n, Ederer, Gin\'{e}, Schmalz},
	colorlinks=true,
	allcolors=blue,
}
%\usepackage[monochrome]{color}
\urlstyle{same}
\usepackage{booktabs}

\usepackage{titlesec}
\setcounter{secnumdepth}{4}
\titleformat{\paragraph}
{\normalfont\normalsize\bfseries}{\theparagraph}{1em}{}
\titlespacing*{\paragraph}
{0pt}{3.25ex plus 1ex minus .2ex}{1.5ex plus .2ex}
% Text markings to edit
\newcommand\red[1]{{\color{red} \textbf{[#1]}}}

\title{\vspace{-2.5cm} Innovation: The Bright Side of Common Ownership?\thanks{We are grateful to Daron Acemoglu, Nick Bloom, Harald Hau, Barry Nalebuff, Luke Taylor, Raffaella Sadun, Regina Seibel, Fiona Scott Morton, John Van Reenen, \#EconTwitter and seminar audiences at Berlin HSG, Humboldt University, NBER PIE, IESE, MadBar Workshop, Michigan, and Yale for helpful comments. We thank Nick Bloom, John Van Reenen, and, in particular, Brian Lucking for generously sharing their data and Pedro Mart\'inez Bruera for outstanding research assistance. Ant\'on acknowledges the financial support of the Department of Economy and Knowledge of the Generalitat de Catalunya (Ref. 2014 SGR 1496), and of the Ministry of Science, Innovation, and Universities (Ref. PGC2018-097335-A-I00). Schmalz acknowledges funding from Deutsche Forschungsgemeinschaft under Germany's Excellence Strategy – EXC 2126/1–390838866.}}
\author{Miguel Ant\'{o}n\thanks{IESE Business School, \href{mailto:manton@iese.edu}{manton@iese.edu}} \and Florian Ederer\thanks{Yale School of Management and Cowles Foundation, \href{mailto:florian.ederer@yale.edu}{florian.ederer@yale.edu}} \and Mireia Gin\'{e}\thanks{IESE Business School, ECGI, CEPR, and WRDS, \href{mailto:mgine@iese.edu}{mgine@iese.edu}} \and Martin Schmalz\thanks{University of Oxford Sa\"{i}d Business School, CEPR, ECGI, CESifo, and C-SEB, \href{mailto:martin.schmalz@sbs.ox.ac.uk}{martin.schmalz@sbs.ox.ac.uk}}} 
\date{\today\vspace{-1cm}}

\begin{document}
\maketitle
\thispagestyle{empty}


\begin{abstract}
\noindent A firm has inefficiently low incentives to innovate when other firms benefit from its innovative activity and the innovating firm does not capture the full surplus of its innovations. We provide conditions under which common ownership of firms mitigates this impediment to corporate innovation. Common ownership increases innovation when technological spillovers are sufficiently large relative to product market spillovers. Otherwise, the business stealing effect of innovation dominates and common ownership reduces innovation. Empirically, product market spillovers (as measured by Jaffe/Mahalanobis proximity in product market space) decrease the effect of common ownership on innovation inputs and outputs, whereas technology spillovers (proximity in patent space) increase the effect. The sign and magnitude of the relationship between common ownership and corporate innovation varies considerably across the universe of firms depending on the relative strength of product market and technology spillovers. When product market spillovers are relatively large, an increase from the 25th to the 75th percentile of common ownership is associated with a \emph{decrease} of -8.4\% in citation-weighted patents. But when technology spillovers are relatively large, the same increase in common ownership is associated with an \emph{increase} of +12.5\% in citation-weighted patents. Our results inform the debate about the welfare effects of increasing common ownership among U.S.~corporations. 

% \noindent A firm has inefficiently low incentives to innovate when other firms benefit from its innovative activity and the innovating firm does not capture the full surplus of its innovations. We theoretically and empirically show under which conditions common ownership of firms can mitigate this impediment to corporate innovation. Common ownership increases innovation when technological spillovers (as measured by firms’ distance to each other in technology space) are sufficiently large relative to product market spillovers (as measured by distance in product market space). Otherwise, the business stealing effect of innovation dominates and common ownership reduces innovation. As a result, common ownership can have a pro- or anticompetitive effect. Our results inform the debate about the welfare effects of the increase of common ownership among U.S. corporations.

% \noindent We model an industry in which firms can choose both cost-reducing R\&D and their product market strategy. The model yields general conditions under which common ownership of the industry competitors can improve unilateral incentives to innovate. We test and find empirical support for the nuanced model predictions. Specifically, common ownership increases innovation when the degree of technological spillovers is high and the number of firms is small. Otherwise, common ownership has a weakly negative effect on innovation. Because a positive effect on R\&D is necessary for common ownership to improve welfare, our findings contribute to understanding the welfare consequences of the increase in common ownership over the past two decades.
\end{abstract}
JEL Codes: O31, L20, L40
% 	O31	Innovation and Invention: Processes and Incentives
%   L2	Firm Objectives, Organization, and Behavior
%   L20	General
%   L4	Antitrust Issues and Policies
%   L40	General

\newpage

\setcounter{page}{1}

\section{Introduction}
\label{sec:intro}

Two secular trends have recently led to a spirited discussion among academics and policy makers regarding the competitiveness of the U.S.~economy. First, increasing levels of product market concentration, as measured at the national industry level, have been accompanied by increasing profitability, a decline of the labor share of income, rising inequality, declining business dynamism, and, perhaps most importantly, declining innovation.\footnote{\citet{cea2016benefits} provides an early overview of these trends. See \cite{philippon2016investment}, \cite{gutierrez2017declining}, \citet{deloecker2020rise}, and \cite{akcigit2021ten} for a formal quantification and analysis of their macroeconomic implications.}
% For prominent press coverage of the topic, see e.g. ``Too much of a good thing: Profits are too high. America needs a giant dose of competition,'' \emph{The Economist}, March 26, 2016 and ``The problem with profits: Big firms in the United States have never had it so good. Time for more competition,'' \emph{The Economist}, March 26, 2016.
% What else should we add here, substantively and strategically? Philippon stuff for example? 
Second, in addition to rising product market concentration and declining innovation, common ownership has also increased: firms are increasingly commonly owned by a decreasing number of institutional investors.\footnote{\citet{backus2021common} provide a recent comprehensive analysis of common ownership of the largest U.S. corporations. See \cite{azar2012new} for an earlier documentation of this trend and \cite{schmalz2018,schmalz2021} for reviews of the theoretical and empirical literature on common ownership.} For example, Softbank's Vision Fund recently attracted the attention of a number of competition authorities by acquiring large stakes in rivals in the ride-hailing industry and exerting its influence to effectuate a lessening of competition in an alleged attempt to ``dominate ride-hailing'' \citep{economist2018ride}. As a result, competition authorities have begun investigations to study the competitive effects of common ownership of industry competitors by mutual funds, hedge funds, and other types of investment vehicles (e.g., Berkshire Hathaway) that pool resources from a large number of investors but concentrate control over portfolio firms in a few asset management firms.\footnote{Mentions of the concerns and investigations by competition authorities and international institutions include, among many others, \cite{oecd2017common}, \cite{eucomp2017decision}, \cite{ftc2018hearing}, \cite{vestager2018competition}, and \cite{bloomberg2020private}.} 
%Add a reference to the Bloomberg on FTC HSR INDEX FUNDS and CCI PRIVATE EQUITY 
Although much attention has focused on the empirical investigation of anticompetitive effects of common ownership, much less work has been devoted to its procompetitive and potentially welfare-enhancing role.

In this paper we investigate, both theoretically and empirically, how corporate innovation depends common ownership. In our model, the sign and the magnitude of the common ownership effect on corporate innovation vary with the relative importance of technological spillovers and business stealing repercussions of innovative activity. We show empirically that the sign and magnitude of the relationship between common ownership and innovation varies considerably across the universe of publicly listed U.S.~corporations, in line with the theory's predictions. 

We begin our analysis by introducing common ownership in a canonical model of (process) innovation and strategic competition with both technology and product market spillovers between firms. Our model allows for product differentiation, technology spillovers, and common ownership to vary across all firm pairs. This permits us to study common ownership links between firms across the entire economy rather than just in a single industry. In the presence of technological spillovers, innovation in one firm not only generates benefits in the firm that produced the innovation but also in technologically related firms. This surplus appropriability problem leads to inefficiently low ex-ante incentives to innovate \citep{bolton1999,jones2000too,arora2021knowledge}. Common ownership of technologically related firms mitigates this problem to the extent that firms act in the interest of these common owners. Common ownership can even render innovative activity profitable that would be unprofitable if it only benefited the innovating firm itself. %When innovation lowers marginal costs in the industry so much as to increase industry output, common ownership may even increase welfare. 
Prior literature has suggested such beneficial knowledge transfers predominantly in the context of private firms \citep{lindsey2008,eldar2020common,gonzalez2020exchanges} or among investors \citep{stein2008conversations,botelho2018here}, whereas we focus on public firms.

However, there is a second dimension affecting the firm's innovation decisions: the interaction between innovation and product market competition. Innovations resulting from R\&D expenditures naturally lead to the innovator stealing market share and profits from firms competing in the same or related product markets \citep{bloom2013}. When the competitors are predominantly owned by separate groups of shareholders, this procompetitive effect of innovation is desirable for the innovating company's shareholders. But when the same shareholders own both the innovator and its product market competitors, such business stealing is less desirable.\footnote{We abstract away from the potential role of common \emph{debtholders} in inducing reduced competition, which is the focus of empirical work by \citet{saidi2021bank}.} Hence, common ownership can reduce the incentives to innovate when the business stealing effect is stronger than the aforementioned technological spillover effect.\footnote{Consistent with this idea, \cite{gonzalez2020exchanges} shows that technological spillovers among companies sharing common VCs are more substantial between portfolio companies that are not in direct competition for the VCs' resources because different funds finance them.} Our theoretical framework combines both of these effects and provides conditions that determine which one of them dominates.

\begin{table}[!htb]
\centering
\medskip
\resizebox{14cm}{!}{
\begin{tikzpicture}
  \path
    (0,1) coordinate (IBM) node[above, inner sep=0] {
    \begin{tabular}{>{\raggedright\arraybackslash}p{4.25cm}>{\raggedleft\arraybackslash}p{1cm}}
    \textit{\textbf{IBM}}        & \textit{[\%]} \\
    \cmidrule(r){1-2} 
    Berkshire Hathaway  &      8.35 \\
    Vanguard            &      6.06 \\
    State Street        &      5.12 \\
    BlackRock           &      5.06 \\
    State Farm          &      1.72 \\
    BNY Mellon          &      1.46 \\
    Fidelity            &      1.29 \\
    Northern Trust      &      1.14 \\
    Norges Bank         &      0.94 \\
    Geode Capital       &      0.75 \\
    \end{tabular}
    }
    (0,-1) coordinate (Motorola) node[below, inner sep=0] {
    \begin{tabular}{>{\raggedright\arraybackslash}p{4.25cm}>{\raggedleft\arraybackslash}p{1cm}}
    \textit{\textbf{Motorola Solutions}} & \textit{[\%]} \\
    \cmidrule(r){1-2} 
    ValueAct            &      10.11 \\
    BlackRock           &      8.67  \\
    Capital Research    &      7.93  \\
    Orbis               &      7.61  \\
    Vanguard            &      5.31  \\
    Parnassus Investments &      4.97 \\
    State Street        &      3.83  \\
    Metropolitan West   &      2.26  \\
    Janus Capital       &      2.09  \\
    Neuberger Berman    &      2.06  \\
    \end{tabular}
    }
    (9,1) coordinate (Intel) node[above, inner sep=0] {
    \begin{tabular}{>{\raggedright\arraybackslash}p{4.25cm}>{\raggedleft\arraybackslash}p{1cm}}
    \textit{\textbf{Intel}} & \textit{[\%]} \\
    \cmidrule(r){1-2}
    BlackRock           &      6.14  \\
    Vanguard            &      6.00  \\
    Capital Research    &      5.56  \\
    State Street        &      3.98  \\
    Wellington          &      2.18  \\
    Northern Trust      &      1.26  \\
    UBS                 &      1.10  \\
    Harris Associates   &      1.09  \\
    BNY Mellon          &      1.01  \\
    Norges Bank         &      0.96  \\
    \end{tabular}
    }
    (9,-1) coordinate (Apple) node[below, inner sep=0] {
    \begin{tabular}{>{\raggedright\arraybackslash}p{4.25cm}>{\raggedleft\arraybackslash}p{1cm}}
    \textit{\textbf{Apple}} & \textit{[\%]} \\
    \cmidrule(r){1-2} 
    Vanguard            &      5.79  \\
    BlackRock           &      5.65  \\
    State Street        &      3.90  \\
    Fidelity            &      2.79  \\
    Northern Trust      &      1.27  \\
    BNY Mellon          &      1.22  \\
    T.~Rowe Price       &      0.90  \\
    Norges Bank         &      0.86  \\
    Invesco             &      0.85  \\
    J.P.~Morgan         &      0.84  \\
    \end{tabular}
    }
    (5,0) coordinate (Middle); %node[right, inner sep=0] {Put here the 3rd table};
\textbf{
    \draw[ultra thick, <->] (-0.15,0.75) -- node[left]{0.46} (-0.15,-0.75); % IBM-Motorola TECH
    \draw[red, ultra thick, <->] (0.15,0.75) --  node[right]{\textcolor{red}{0.01}} (0.15,-0.75); % IBM-Motorola SIC
    \draw[ultra thick, <->] (8.85,0.75) -- node[left]{0.47} (8.85,-0.75); % Intel-Apple TECH
    \draw[red, ultra thick, <->] (9.15,0.75) -- node[right]{\textcolor{red}{0.00}} (9.15,-0.75); % Intel-Apple SIC
    \draw[ultra thick, <->] (3.5,4.15) -- node[above]{0.76} (5.5,4.15); % IBM-Intel TECH
    \draw[red, ultra thick, <->] (3.5,3.85) -- node[below]{\textcolor{red}{0.01}} (5.5,3.85); % IBM-Intel SIC
    \draw[ultra thick, <->] (3.5,-3.85) -- node[above]{0.17} (5.5,-3.85); % Motorola-Apple TECH
    \draw[red, ultra thick, <->] (3.5,-4.15) -- node[below]{\textcolor{red}{0.02}} (5.5,-4.15); % Motorola-Apple SIC
    \draw[ultra thick, <->] (3.15,1.25) node[above, inner sep=2]{\phantom{xxxxx}0.64} to [out=0,in=90](6.25,-1); % IBM-Apple TECH
    \draw[red, ultra thick, <->] (3.15,1) node[below, inner sep=6]{\phantom{xxxxx}\textcolor{red}{0.65}} to [out=0,in=90](6,-1); % IBM-Apple SIC
    \draw[ultra thick, <->] (3.15,-1) node[above, inner sep=6]{\phantom{xxxxx}0.46} to [out=0,in=-90](6,1); % Motorola-Intel TECH
    \draw[red,ultra thick, <->] (3.15,-1.25) node[below, inner sep=2]{\phantom{xxxxx}\textcolor{red}{0.34}} to [out=0,in=-90](6.25,1); % Motorola-Intel SIC
}
\end{tikzpicture}
}
\protect\caption{Ownership holdings and technology \& product market correlations of four high technology firms{\scriptsize{} \newline The table presents the ten largest owners and their ownership holdings in 2015 of the four technology companies (IBM, Intel, Motorola, and Apple) discussed in the text. The arrows between each of the company pairs depict the technology (black) and product market (\textcolor{red}{red}) correlations reported in \cite{bloom2013}.}}
\label{tab:owners}
\end{table}

Including both dimensions is of first-order importance for understanding the overall effect of common ownership on innovation not just in the theory but also in our empirical implementation. To illustrate, Table~\ref{tab:owners} reports the ownership shares of four technology firms (IBM, Intel, Motorola, and Apple) that are technologically related but compete to a varying extent in the same product markets. First, the four companies are closely technologically related over the sample period. The technological distances, as measured by firms' patent issuances across different patent classes in \cite{bloom2013}, between IBM-Intel, IBM-Motorola, and IBM-Apple are 0.76, 0.46, and 0.64, respectively---much larger than the sample average of 0.038. Product market proximity is more hetergeneous across these firm pairs. Whereas IBM is close to Apple in product market space (product market distance of 0.65, compared to the sample average of 0.015), IBM is not close to Intel and Motorola (product market distances of 0.01). 

As shown in Table~\ref{tab:owners}, these four firms also have a significant degree of common ownership, particularly toward the end of our sample period in 2015. BlackRock, Vanguard, and State Street are all represented among the top owners of each of the four companies. However, the large concentrated blockholdings by Berkshire Hathaway and ValueAct illustrate that common ownership interests are heterogeneous and asymmetric across firms. Therefore, the degree of common ownership between firms may differentially affect their innovation decisions as a function of the firms' technological and product market proximities. Our central theoretical prediction is that the effect of common ownership on innovation depends on firms' relative distances in technology and product market space and can vary both in terms of sign and magnitude.

Whether the theoretical predictions about the relationship between common ownership and innovation are helpful in organizing the data is a question that requires more than just anecdotal evidence. We empirically show that both effects exist and that they lead to substantial heterogeneity of the relationship between common ownership and innovation across firms. To do so, we use the methodology pioneered by \citet{bloom2013} and \citet{lucking2018spillovers} to measure technology and product market spillovers.
%\footnote{Both effects can lead to improvements in firm value. We therefore do not examine how stock market reactions relate to changes in common ownership as in \citet{boller2019testing} with their effect of innovation.}
We combine these data with information about the ownership of firms, in particular to which extent the largest owners of one firm also hold shares in other firms using the ``kappa'' measure advocated by \cite{backus2021common}. Using panel regressions we find an ambiguous relationship \emph{on average} between common ownership and corporate innovation as measured by innovation inputs (scaled R\&D expenditures) and innovation outputs (citation-weighted number of patents and total stock market value of patents). However, throughout all of our specifications, innovation is more positively related to common ownership when technological spillovers are higher, whereas more common ownership is associated with less innovation when product market spillovers are greater. As a result (and as predicted by our theory), common ownership and corporate innovation are positively related when technology spillovers are large relative to product market business stealing incentives and are negatively related otherwise. Finally, the relationship between common ownership and innovation varies considerably across the publicly listed firms depending on the relative strength of product market and technology spillovers. For example, when product market spillovers are relatively large, an increase from the 25th to the 75th percentile of common ownership is associated with a \emph{decrease} of -8.4\% in citation-weighted patents. In contrast, when technology spillovers are relatively large, the same increase in common ownership is associated with an \emph{increase} of +12.5\% in citation-weighted patents.

%Both innovation inputs and outputs
%Common ownership measured as kappas

% To provide a tighter link between the model which assumes that ownership is exogenous, and empirical tests, in which ownership may be endogenous due to unmodeled equilibrium considerations, we employ a difference-in-differences strategy that allows us to shock common ownership of the treated firm without affecting its ownership structure. Specifically, we investigate how S\&P 500 firm's innovation inputs and outputs change after a product market competitor is added to the index. %TBC What the heck do we do? Add product market competitors or add technologically related firms?

% \subsection{Related Literature}
% Innovation is the primary driver of economic growth and research and development (R\&D) has therefore been a major topic in a variety of literatures on corporate innovation and corporate strategy, and corporate governance for many decades. 

Given that incentives to compete are tightly linked to incentives to innovate \citep{daspremont1988,hoskisson2002conflicting,aghion2005,bloom2013}, our paper lies at the intersection of corporate innovation, corporate strategy, and corporate governance. The extant literature focuses on the potential benefits of cooperative R\&D, on how innovation is affected by mergers, or how it relates to institutional ownership.\footnote{For the interplay between competition and innovation see, for example, \citet{brander1983}, \citet{spence1984}, \citet{katz1986}, \citet{daspremont1988}, \citet{grossman1991}, \citet{kamien1992}, \citet{suzumura1992}, \citet{aghion1992}, and \citet{leahy1997}. For comprehensive reviews of the literature see \citet{jones2005} and \citet{gilbert2006}.} One of this literature's primary objectives is to examine the underinvestment of R\&D and the welfare effects of moving from a noncooperative to a cooperative regime in R\&D. For example, \cite{kamien1992} identify conditions under which a cartelized research joint venture (RJV) is optimal. \cite{leahy1997} show that R\&D cooperation leads to more output, innovation, and welfare when spillovers are positive. We adopt these canonical models of innovation and product market competition and re-examine their conclusions in light of the fact that firms with different names do not necessarily have disjoint sets of investors.
%Shall we mention all the shitty and less related papers on common ownership and innovation?

The most closely related paper to our own analysis is \cite{bloom2013} who study the effect of product market and technology spillovers on innovation and provide economy-wide empirical evidence for the importance of both effects, but without considering the role of common ownership. They estimate the extent of spillovers in a panel of U.S.~firms from 1981 to 2001 and find that gross social returns to R\&D are at least twice as high as the private returns. Their results imply that the internalization of those technological spillovers is a matter of first-order welfare importance.\footnote{Their approach builds on prior work by \citet{jaffe1988} who assigns firms to technology and product market space but does not examine the distance between firms in both these spaces. Similarly, \citet{branstetter2002} empirically examine the effects of technology closeness and product market overlap on patenting in Japanese research consortia. \citet{lucking2018spillovers} extend the results of \cite{bloom2013} to later time periods.} We demonstrate that common ownership can internalize product market and technology externalities between the firms and thereby significantly affect the level and heterogeneity of corporate innovation. Our paper is also related to \cite{lopez2016cross} who theoretically study the effect of (symmetric and identical) common ownership on innovation of several symmetric industry competitors.\footnote{\cite{stenbacka2020overlapping} theoretically study how common ownership affects the product innovation decisions under duopoly and distinguishes between input spillovers and output spillovers.} In their model, all firms are symmetric, compete in the same industry, and produce undifferentiated products. Technology spillovers and common ownership shares are identical between them. In contrast, our model allows for common ownership of firms in the entire economy, including potentially in separate industries. To reflect that greater scope, we allow for product differentiation, technology spillovers, and common ownership to vary across firms. These generalizations are crucial to predict and understand the variation of the effect of common ownership on innovation found in the data. % That's a fine assumption for tractability reasons in their case which is just theoretically studying a single industry, but less applicable for the economy-wide empirical implementation that we are pursuing.
%
% product market spillovers fixed (they do consider comparative statics with technology spillovers), but both product market and technology spillovers are assumed to be the same. That's a fine assumption for tractability reasons in their case which is just theoretically studying a single industry, but less applicable for the economy-wide empirical implementation that we are pursuing. But you're right, we should make that much clearer in the draft.

Several recent contributions provide distinct empirical investigations of the relationship between various measures of common ownership and innovation. Our analysis differs from all of them in that it employs theoretical model that guides our empirical design and informs the interpretation of our results. Our analysis may help explain the substantial variation in the sign and magnitude of common ownership effects on innovation across these different papers. \cite{li2020venture} study common venture capital ownership of pharmaceutical startups and find evidence suggesting that common ownership improves innovation efficiency. In contrast to their work, we focus on a broad sample of public firms. \cite{he2017product} examine the question of whether common \emph{blockholders} have an effect on corporate innovation on average. In contrast, we study the entire institutional ownership structure of the firm and examine whether the degree of technology spillovers and product market spillovers differentially affect the relationship between common ownership and innovation. \cite{kostovetsky2020common} show that increases in shared institutional ownership caused by index additions are followed by more citations of the patents of the firm that was added to the index. \cite{borochin2020common} provide evidence that the sign of the relationship between patent output, non-self citations, and common ownership depends on the type of institutional owner creating the common ownership link. \cite{chiao2020corporate} argue that common ownership (as measured by industry-level MHHID in a sample ending in 2008) is on average negatively related to patent grants, citations, and R\&D expenditures. They also find that common ownership is negatively associated with the likelihood that firms are involved in patent litigation and positively related to the speed of the settlement of lawsuits between commonly owned firms. The empirical evidence of \cite{geng2016technological} suggests that common ownership can mitigate hold-up problems between firms owning complementary patent portfolios. However, the strength of the effect depends on the type of institutional owner. Finally, using a pan-European sample, \cite{gibbon2021rising} find that common ownership (as measured by MHHI delta calculated at the three-digit industry and country level) is related to citation-weighted patents in high technology industries, whereas common ownership is related to markups in low technology industries.


%\footnote{Less closely related is a literature in management and finance on the effects of acquisitions on innovation and vice versa \citep{hitt1991effects,hitt1996market,bena2014corporate}, as well as a literature on the link between innovation, firm creation, and asset prices \citep{bena2016heterogeneous}. Interestingly, none of the papers in the finance literature studies the role of shareholder diversification across competitors on innovation.}
%However, ownership structure has long been linked to how firms choose to compete and innovate. \cite{lindsey2008} documents that strategic alliances are more frequent among companies sharing a common venture capitalist, but only when at least one of the portfolio firms is private and with the effect being particularly strong for industry competitors.\footnote{In a related vein, \citet{aghion2013innovation} find that greater institutional ownership is (causally) related to more innovation and \citet{brav2014shareholder} show that firms targeted by hedge fund activists experience a tightening in R\&D expenditures, but an improvement in innovation efficiency following the intervention.} However, in theory, common ownership can increase or decrease firms' incentives to innovate, depending on two countervailing forces which stem from two distinct types of spillover effects of R\&D.
%technology (or knowledge) spillovers, which may increase the productivity of other firms that operate in similar technology areas; second, product market rivalry effect of R\&D. Whereas technology spillovers are beneficial to other firms, R\&D by product market rivals has a negative effect on a firm's profits due to business stealing.


%In contrast, contributions to the literature on overlapping ownership starting with \cite{bresnahan1986} and \cite{reynolds1986} have mainly focused on its anticompetitive effects. \citet{azar2015anti, azar2016bank} provide evidence that common ownership causes higher product prices in the airline and banking industries, respectively. \citet{philippon2016investment} show that firms owned by quasi-indexers tend to underinvest relative to investment opportunities in a broad panel of US firms. \citet{xie2018institutional} document that pharmaceutical brands pay commonly owned generic manufacturers for delayed entry. \citet{anton2016} document the secular increase of common ownership among publicly listed US firms and show that it is related to performance-insensitive compensation, thus providing evidence of the transmission mechanism of common institutional ownership on managers' incentives. The potentially beneficial effects of common ownership are the primary focus of \citet{lopez2016} who theoretically analyze under what conditions common ownership is welfare-enhancing and what the socially optimal level of common ownership is. Our paper extends existing theoretical results and provides the first empirical analysis of the countervailing product and technology spillover effects of common ownership on corporate innovation.

The remainder of this paper is organized as follows. Section~\ref{sec:model} presents the theoretical framework that guides the empirical analysis. Section~\ref{sec:data} describes the data. The empirical results are presented and discussed in Section~\ref{sec:results}. Section~\ref{sec:conclusion} concludes.

%This paper contributes to a large literature on the drivers of corporate innovation and a fast growing literature on common ownership. With respect to the theoretical part of our contribution, \citet{lopez2016} is the most closely related paper. The authors derive conditions under which common ownership can increase innovation, output, and welfare in a Cournot setup in which all firms produce a homogeneous product. Our relative contribution is an extension of their model to allow for product differentiation to capture differential business stealing incentives between firms based on their proximity in product market space.\footnote{We also show that our results extend to the case of Bertrand competition.} We use this model to derive two important predictions which relate the sign of the common ownership effect on innovation to technological and product market spillovers.\footnote{Relative to standard models of competition under common ownership (e.g., \citet{salop2000common}), our model allows for innovation with technological spillovers in addition to quantity choices. In our model, an investment in R\&D not only reduces the firm's own costs, but also spills over to closely technologically related firms.}
%
%
%Our main contribution, however, is an empirical investigation of the effect of common ownership in the presence of both technology and product market spillovers.
%
%%\citet{utterback1994mastering} 


\section{Theoretical Framework} 
\label{sec:model}

\subsection{Setup}

We analyze the role of common ownership and its interplay with product market and technological spillovers in the canonical model of innovation and product market competition pioneered by \cite{daspremont1988}. By doing so, we also extend the model of \cite{bloom2013}, which studies the effect of product market and technology spillovers on innovation, to allow for overlapping ownership between firms. Our theoretical setup is also related to the model of \cite{lopez2016cross} which studies the interplay between innovation and common ownership, but we allow for both product market and technology spillovers as well as common ownership weights to differ across firms. 

% That's a fine assumption for tractability reasons in their case which is just theoretically studying a single industry, but less applicable for the economy-wide empirical implementation that we are pursuing.
%
% product market spillovers fixed (they do consider comparative statics with technology spillovers), but both product market and technology spillovers are assumed to be the same. That's a fine assumption for tractability reasons in their case which is just theoretically studying a single industry, but less applicable for the economy-wide empirical implementation that we are pursuing. But you're right, we should make that much clearer in the draft.

Firms' innovation choices, product quantities, prices, and profits are endogenously co-determined by the degree of common ownership as well as product market and technological spillovers. In line with the existing literature on common ownership, we assume that ownership is exogenous.

\subsubsection{Product Market Competition}

Consider an economy with $n$ firms that each produce a single differentiated product. There are no industries per se, but all firms compete with each other depending on how closely related their products are. In our model, the welfare-enhancing effects of common ownership are due to economy-wide horizontal and vertical externalities that arise from technology spillovers. Although strictly speaking we present a partial equilibrium model, all of our insights regarding the impact of common ownership under different technology and product market spillovers also hold in a general equilibrium model.\footnote{For example, \cite{pellegrino2019product} and \cite{ederer2021welfare} model and estimate a general equilibrium hedonic linear demand system in which all the $n$ firms in the economy compete with each other and the investors (or managers) controlling the firms' operations consume an outside good (e.g., leisure).}

Following \cite{singhvives1984} and \cite{hackner2000}, we derive demand from the behavior of a representative consumer with the following quadratic utility function:
\begin{equation}\label{eq:representative}
U(\mathbf{q}) = A \sum_{i=1}^{n} q_{i} - \frac{1}{2}\left(a_{ii} \sum_{i=1}^{n} q_{i}^{2} + 2 \sum_{i \neq j} a_{ij} q_{i} q_{j}\right)
\end{equation}
where $q_i$ is the quantity of product $i$, $\mathbf{q}=(q_{1},...,q_{n})$ is the vector of all quantities, $A > 0$ represents overall product quality, $a_{ii} > 0$ measures the concavity of the utility function, and $a_{ij}$ represents the degree of substitutability between two differentiated products $i$ and $j$. $a_{ii} > a_{ij} \geq 0$ ensures that the products are (imperfect) substitutes. Without loss of generality and to simplify notation, we set $a_{ii} = 1$. The higher the value of $a_{ij}$, the more alike the products are. The resulting consumer maximization problem yields linear demand for each product $i$, such that the firms face symmetric inverse demand functions given by
\begin{align}\label{eq:market_demand}
p_{i}(\mathbf{q}) = A - q_{i} - \sum_{j \neq i}^{n} a_{ij} q_{j},
\end{align}
where $i = 1, 2, ..., n$. Because $1 > a_{ij} \geq 0$, a firm's quantity $q_{i}$ has a greater impact on the price $p_{i}$ for its own product than the quantity of any other firm $q_{j}$.\footnote{In the main part of the paper we focus on the Cournot competition case where quantity choices are strategic substitutes. However, our results for Bertrand competition (see Appendix) where prices are strategic complements are essentially identical. Although we assume linear demands, the main results of our model generalize to nonlinear demand functions.} The parameter $a_{ij}$ measures product homogeneity or product market spillovers. Given the symmetry of the empirical measure of product market spillovers \citep{bloom2013} that we describe in Section~\ref{sec:data}, we assume that this parameter is symmetric between firm $i$ and $j$, $a_{ij} = a_{ji}$. If $a_{ij}$ is small, the products of firm $i$ and $j$ are quite distinct and thus expanding output $q_{i}$ (or lowering price $p_{i}$) does not steal much market share from the competing firm $j$. On the other hand, if $a_{ij}$ is large the product varieties produced by the firms are quite similar and thus business stealing is more pronounced.

\subsubsection{Innovation}

Following the extant theoretical literature on innovation \citep{daspremont1988,kamien1992,leahy1997,lopez2016cross} we model corporate innovation as decreasing marginal cost. This is just a particular modeling choice that ensures tractability. One could also model innovation as increasing product quality which would yield qualitatively similar results. As in much of the existing innovation literature we only focus on the intensity of innovation, but do not consider the direction of innovation which is the focus of \cite{letina2016road}, \cite{bryan2017direction}, and \cite{callander2021novelty}. Common ownership may also influence which innovation projects firms choose to pursue. 

Firm $i$ has a marginal cost of $c_{i}$ given by
\begin{align}
c_{i} = \bar{c} - x_{i} - \sum_{j \neq i}^{n} \beta_{ij} x_{j}.
\end{align}
Firm $i$ can lower its marginal cost from $\bar{c}$ by investing in innovation $x_{i}$ at a cost $\frac{\gamma}{2} x_{i}^2$. A firm's marginal costs are also reduced by the innovative investments of other firms $x_{j}$, to the extent that these investments benefit firm $i$ because of technological spillovers captured by $0 \leq \beta_{ij} < 1$. This means that a firm $i$'s investment in innovation reduces its own marginal cost $c_{i}$ and to a lesser extent may also reduce the marginal cost $c_{j}$ of firm $j$. Given the construction of the empirical measure of technological spillovers \citep{bloom2013} we assume that this parameter is symmetric, $\beta_{ij} = \beta_{ji}$. These technological spillovers are not confined within the same industry or even just to firms that produce relatively similar substitute products. Innovation benefits can spill over to technologically related firms (i.e., $\beta_{ij} > 0$) that produce goods that are entirely unrelated in terms of product market competition (i.e., $a_{ij} = 0$). The example mentioned in the introduction of IBM and its relationship to Intel and Motorola, which are close in technology space but not in product market space, fits this case quite well.  

The profits of firm $i$ are given by
\begin{align}
\pi_{i} = q_{i}\left[A - q_{i} -  \sum_{j \neq i}^{n} a_{ij} q_{j} - \left(\bar{c} - x_{i} - \sum_{j \neq i}^{n} \beta_{ij} x_{j}\right) \right] - \frac{\gamma}{2} x_{i}^2.
%\pi_{i} = (p_{i}-c)(A-bp_{i}+ap_{j})\\
\end{align}
Firms choose quantities $q_{i}$ and innovation levels $x_{i}$ simultaneously. We obtain qualitatively similar results when firms invest in innovation before choosing quantities (or prices).

\subsubsection{Owners}

There are $n$ owners which share the same index as the $n$ firms. Each owner $i$ owns a stake in firm $i$ as well as shares in other firms denoted by $j \neq i$. We assume firms act in their largest owners' interest. Our model nests the special case in which firms maximize their own profits.\footnote{Assuming shareholders agree on own-firm value maximization has no theoretical basis when firms are not price takers and shareholders have interests outside the firm \cite{hart1979shareholder}. Furthermore, firms acting in their shareholders' portfolio interest is also a better description of how firms behave and how managers are incentivized \citep{anton2021incentives}.}$^,$\footnote{Aside from a literal interpretation, this assumption can be understood as a metaphor for an explicit or implicit coalition of shareholders that jointly hold an effective majority of the voting stocks. \cite{olson2017fracking} and \cite{shekita2020interventions} discuss examples of explicit coalitions. \cite{moskalev2020} shows conditions under which shareholders with similar portfolios will optimally vote the same way, and therefore will be regarded as an implicit coalition or a single block by managers.}

% To simplify the exposition, we assume that these owners are symmetric such that owner $i$ owns a majority stake in firm $i$ and a residual symmetric share in all other firms. 
Specifically, we follow the common ownership literature since \cite{rotemberg1984financial} and assume that firms maximize a weighted average of their shareholders' portfolio profits. \cite{azar2012new} and \cite{backus2021common} show that firm $i$'s maximization problem can be restated as:
\begin{align}\label{eq:objective_function}
\phi_{i} = \pi_{i} + \sum_{j \neq i} \kappa_{ij} \pi_{j}  
\end{align}
where $\kappa_{ij}$ is the weight that firm $i$ places on the profits $\pi_j$ of firm $j$. Its exact value depends on the type of ownership and corresponds to what \citet{edgeworth1881} termed the ``coefficient of effective sympathy among firms.'' In fact, even before the common ownership literature there is a long tradition in economics of weighting shareholder interests in the objective function of the firm, including \cite{dreze1974investment} and \cite{grossman1979theory}. We assume that the profit weight $\kappa_{ij}$ is between 0 (separate ownership) and 1 (perfectly common ownership). In contrast to $a_{ij}$ and $\beta_{ij}$, we do not assume that $\kappa_{ij}$ is symmetric between any firm pair $i$ and $j$, that is $\kappa_{ij} \neq \kappa_{ji}$ in general. 

We use the $\kappa$ notation of \cite{backus2021common} which is equivalent to $\lambda$ in \cite{azar2012new}, \cite{lopez2016cross}, and \cite{azar2020oligopoly}. Values of $\kappa$ exceeding 1 are possible, but they lead to owners placing more weight on their competitors' profits than on their own profits. This would make it possible for common ownership to create incentives for the ``tunneling'' of profits from one firm to another \citep{johnson2000tunneling}. By maximizing equation (\ref{eq:objective_function}), the owner essentially maximizes a weighted average of her own firm as well as other firms' profits that she owns. The particular objective function given in equation (\ref{eq:objective_function}) is a normalization. Firms do not maximize a sum that is larger than the entire economy. 
% In particular, \citet{lopez2016cross} consider two types of minority shareholdings: when investors acquire firms' shares (\emph{common ownership}) with silent financial interest or proportional control, and when firms acquire other firms' shares (\emph{cross-ownership}). In both cases, they show that when the stakes are symmetric, firm/owner $i$'s problem is to maximize the objective function given in equation (\ref{eq:objective_function}).

\subsection{Analysis and Comparative Statics}

We now analyze the differential impact that common ownership has on corporate innovation that depends on both product market and technological spillovers. Firm $i$'s first-order conditions with respect to quantity $q_{i}$ and innovation $x_{i}$ can be rearranged to yield the following best-response functions
\begin{align}
q_{i} &= \frac{1}{2} \left[ A - \left( \bar{c} - x_{i} - \sum_{j \neq i}^{n} \beta_{ij} x_{j} \right) - \sum_{j \neq i}^{n} a_{ij} q_{j} - \underbrace{\sum_{j \neq i}^{n} \kappa_{ij} a_{ji} q_{j}}_{\textrm{CO $\times$ product market spillovers}}  \right] \label{BR_quantity} \\
x_{i} &= \frac{1}{\gamma} \left(q_{i} + \underbrace{\sum_{j \neq i}^{n} \kappa_{ij} \beta_{ji} q_{j}}_{\textrm{CO $\times$ technology spillovers} } \right) 
\label{BR_innovation}
\end{align}
where given our symmetry assumptions $a_{ij} = a_{ji}$ and $\beta_{ij} = \beta_{ji}$.

Firm innovation $x_{i}$ is directly proportional to firm quantity $q_{i}$ such that any increase in quantity $q_{i}$ will also increase innovation $x_{i}$. These first-order conditions illustrate the driving forces of our model. When common ownership $\kappa_{ij}$ increases, this has two distinct effects on firm $i$'s first-order conditions. 

First, in equation (\ref{BR_quantity}) an increase in $\kappa_{ij}$ reduces $q_{i}$ through the interaction of common ownership and product market spillovers (i.e., the term labeled ``CO $\times$ product market spillovers'') and thereby reduces innovation $x_{i}$ in equation (\ref{BR_innovation}). This is the anticompetitive effect of common ownership arising from product market spillovers. Effectively, increasing innovation $x_{i}$ causes firm $i$ to steal business from any firm $j$ that is selling a substitute product. This well-known business stealing effect of innovation will be larger the greater the product homogeneity (also known as the degree of product market spillovers) $a_{ij}$. The more closely related the products are, the larger will be the negative profit impact of an increase in quantity on other firms. Common ownership exacerbates this negative effect of product market similarity $a_{ij}$ on output and innovation, because common ownership weakens the firm's business stealing incentive. The reason is that when a firm's objective function puts positive weight $\kappa_{ij}$ on other firms' profits $\pi_{j}$, firm $i$ will partly internalize any negative profit repercussions on these other firms by reducing innovation $x_{i}$ and quantity produced $q_{i}$.

Second, in equation (\ref{BR_innovation}) an increase in $\kappa_{ij}$ directly increases innovation. When firm $i$ innovates, it benefits other firms $j \neq i$ by lowering their marginal costs $c_{j}$. This is the procompetitive effect of common ownership arising from technological spillovers (i.e., the term labeled ``CO $\times$ technology spillovers''). The greater the technological proximity $\beta_{ij}$ between the two firms, the larger is this technology spillover effect. This is because firm $j$ which is more closely located in technology space to firm $i$, will benefit more from the firm $i$'s innovation. Common ownership strengthens this technology spillover effect because with a positive weight $\kappa_{ij}$ in its objective function, firm $i$ partly internalizes the positive externality of innovation on other firms $j \neq i$ that it would otherwise ignore.
% \footnote{We model technological spillovers as a positive externality, but the technological proximity measure we use in the empirical implementation does not sign the externality. The basic prediction that common ownership internalizes the innovation externality is general. \cite{geng2016technological} study the specific case of a negative externality of one patent on another and provide evidence suggesting that common ownership mitigates holdup problems.} 
This output-increasing technology spillover effect is still present when the firms have no product market connection ($a_{ij}=0$). In graphical terms, an increase in $\kappa_{ij}$ tilts the innovation reaction function of firm $i$ inward due to the product market spillovers operating through $a_{ji}$, but shifts it outward due to the technology spillovers operating through $\beta_{ji}$.

It is immediately obvious that the effect of common ownership on innovation has an ambiguous sign: it can be either positive or negative depending on the relative strength of the product market and technology spillovers. If $a_{ij} = 0$ (i.e., product market spillovers are absent) any increase in common ownership $\kappa_{ij}$ will raise firm innovation $x_{i}$ due to technological spillovers $\beta_{ij} \geq 0$. Conversely, if $\beta_{ij} = 0$ (i.e., technological spillovers do not exist), any increase in $\kappa_{ij}$ will decrease firm innovation $x_{i}$ due to product market spillovers $a_{ij} \geq 0$.

We can rewrite the system of first order conditions given in equations (\ref{BR_quantity}) and (\ref{BR_innovation}) in the following way
\begin{align*}
    \left(\mathbf{a} + \mathbf{K} \circ \mathbf{a}'\right) \mathbf{q} &= (A-\bar{c}) \cdot \mathbf{1} + \mathbf{B} \mathbf{x} \\
    \left(\mathbf{K} \circ \mathbf{B}'\right) \mathbf{q} &= \gamma \mathbf{x}
\end{align*}
where $\circ$ is the Hadamard (element-by-element) product, $\mathbf{1}$ is an $n \times 1$ vector of ones, $\mathbf{a}$ is the product similarity matrix, $\mathbf{B}$ is the technology spillover matrix, and $\mathbf{K}$ is the common ownership matrix. The matrices $\mathbf{a}$, $\mathbf{B}$, and $\mathbf{K}$ are defined as follows:
\begin{align*}
    \mathbf{a} = \begin{bmatrix}
	1 & a_{12} & \cdots & a_{1n}\\
	a_{21} & 1  & \cdots & a_{2n}\\
	\vdots & \vdots &\ddots &\vdots \\
	a_{n1} & a_{n2}& \cdots& 1
	\end{bmatrix}, \quad
    \mathbf{B} = \begin{bmatrix}
	1 & \beta_{12} & \cdots & \beta_{1n}\\
	\beta_{21} & 1  & \cdots & \beta_{2n}\\
	\vdots & \vdots &\ddots &\vdots \\
	\beta_{n1} & \beta_{n2}& \cdots& 1
	\end{bmatrix}, \quad
    \mathbf{K} = \begin{bmatrix}
	1 & \kappa_{12} & \cdots & \kappa_{1n}\\
	\kappa_{21} & 1  & \cdots & \kappa_{2n}\\
	\vdots & \vdots &\ddots &\vdots \\
	\kappa_{n1} & \kappa_{n2}& \cdots& 1
	\end{bmatrix}
\end{align*}


% Rewriting the first-order conditions in matrix form yields
% \begin{align*}
% 	\underbrace{\begin{bmatrix}
% 	2 & (1+\kappa_{12})a_{12} & \cdots & (1+\kappa_{1n})a_{1n}\\
% 	 (1+\kappa_{21})a_{21} & 2  & \cdots & (1+\kappa_{2n})a_{2n}\\
% 	\vdots & \vdots &\ddots &\vdots \\
% 	(1+\kappa_{n1})a_{n1} &  (1+\kappa_{n2})a_{n2}& \cdots& 2
% 	\end{bmatrix}}_{:=\mathbf{K_a}}
% 	\underbrace{\begin{bmatrix}
% 	q_1\\ q_2 \\ \vdots \\ q_n
% 	\end{bmatrix}}_{:=\mathbf{q}}
% 	&
% 	=(A-\bar{c})\begin{bmatrix}
% 	1\\1\\ \vdots \\ 1
% 	\end{bmatrix}+\underbrace{\begin{bmatrix}
% 	1 & \beta_{12} & \cdots & \beta_{1n}\\
% 	\beta_{21} & 1  & \cdots & \beta_{2n}\\
% 	\vdots & \vdots &\ddots &\vdots \\
% 	\beta_{n1} & \beta_{n2}& \cdots& 1
% 	\end{bmatrix}}_{:=\mathbf{B}}
% 	\underbrace{\begin{bmatrix}
% 	x_1\\x_2\\ \vdots \\ x_n
% 	\end{bmatrix}}_{:=\mathbf{x}}
% 	\\
% 	\underbrace{\begin{bmatrix}
% 	1 & \kappa_{12}\beta_{12} & \cdots & \kappa_{1n}\beta_{1n}\\
% 	\kappa_{21}\beta_{21} & 1  & \cdots & \kappa_{2n}\beta_{2n}\\
% 	\vdots & \vdots &\ddots &\vdots \\
% 	\kappa_{n1}\beta_{n1} & \kappa_{n2}\beta_{n2}& \cdots& 1
% 	\end{bmatrix}}_{:=\mathbf{K_\beta}}
% 	\begin{bmatrix}
% 	q_1\\ q_2 \\ \vdots \\ q_n
% 	\end{bmatrix}
% 	&
% 	=
% 	\gamma \begin{bmatrix}
% 	x_1\\x_2\\ \vdots \\ x_n
% 	\end{bmatrix}
% \end{align*}

Defining $\mathbf{K_{a}} = \mathbf{a} + \mathbf{K} \circ \mathbf{a}'$ and $\mathbf{K_{\beta}} = \mathbf{K} \circ \mathbf{B}'$ and substituting the second system of first-order conditions into the first system yields the vector of equilibrium innovation choices $\mathbf{x^*}$ given by
\begin{align}
	\mathbf{x^*}
	= \begin{bmatrix} x^*_1\\x^*_2\\ \vdots \\ x^*_n \end{bmatrix}
	= (A-\bar{c})
	\left[
	\gamma \mathbf{K_a}\mathbf{K_\beta}^{-1}-\mathbf{B}\right]^{-1} \cdot \mathbf{1} \label{eq:equilibrium_innovation}
\end{align}
where $\mathbf{1}$ is an $n\times 1$ vector of ones.

\begin{proposition}\label{proposition_main}
Common ownership $\kappa_{ij}$ increases equilibrium firm innovation $x^{*}_{i}$ if and only if technological spillovers $\beta_{ij}$ are sufficiently large relative to product market spillovers $a_{ij}$. The effect of $\kappa_{ij}$ on $x_{i}^{*}$ is decreasing in $a_{ij}$, $\frac{\partial^2 x_{i}^{*}}{\partial \kappa_{ij} \partial a_{ij}} < 0$, and increasing in $\beta_{ij}$, $\frac{\partial^2 x_{i}^{*}}{\partial \kappa_{ij} \partial \beta_{ij}} > 0$.
\end{proposition}

Proposition~\ref{proposition_main} shows that without knowledge of product differentiation and technological characteristics common ownership has an ambiguous effect on innovation. This insight may help explain some of the variation in empirical findings to date on the relation between common ownership and corporate innovation \citep{kostovetsky2020common,borochin2020common,chiao2020corporate}. These empirical designs do not make the distinctions that our theoretical framework predicts to be crucial for determining the sign of the effect of common ownership on innovation. Depending on the relative strengths of (i) the business stealing and (ii) the technology spillover effect, common ownership can either decrease or increase corporate innovation. However, our framework also predicts under what conditions common ownership has a negative or a positive effect on innovation. Common ownership should decrease innovation if $a_{ij}$ is sufficiently large relative to $\beta_{ij}$, whereas common ownership should increase innovation if the opposite is the case. In other words, we expect common ownership to decrease (increase) innovation when product market spillovers are sufficiently large (small) and technology spillovers are sufficiently small (large).

In our empirical implementation we follow \cite{bloom2013} and construct measures of firm-specific product market spillovers for $\sum_{j \neq i}^{n} a_{ji} q_{j}$ and of firm-specific technological spillovers for $\sum_{j \neq i}^{n} \beta_{ji} q_{j}$ which we interact with the firm-specific common ownership measure of $\kappa_{ij}$ to estimate the pro- and anticompetitive effects of common ownership due to product market and technology spillovers highlighted in the first-order conditions (\ref{BR_quantity}) and (\ref{BR_innovation}).\footnote{\cite{bloom2013} provide microeconomic foundations for the spillover measures we use. However, neither they nor we address the ``reflection problem'' pointed out by \cite{acemoglu2014}.}

Importantly, once we control for the relative strength of product market and technology spillovers, the sign of the effect of common ownership on innovation is unambiguous and does not depend on whether firms compete in strategic substitutes or strategic complements. This is in contrast to the analysis in \cite{bloom2013} where many of the predictions depend on the form of strategic competition. The reason for this difference is that common ownership has a ``direct effect'' (i.e., directly affecting the objective function) rather than a ``strategic effect'' (i.e., indirectly affecting it through the effect on decisions of other firms) as defined by \cite{fudenberg1984fat}. Hence, the sign of the common ownership effect on innovation does not depend on the sign of the strategic response of other firms.

Our predictions 
%obtained from Corollary \ref{corollary_comparative_statics} and \ref{corollary_cross_derivatives} 
provide theoretical guidance for our empirical analysis. Specifically, they allow us to quantify whether and under what conditions common ownership should increase or decrease innovation and how product market and technology spillovers should affect this relationship.


\section{Data}
\label{sec:data}

In this section we investigate the empirical relationship between common ownership, product market competition, and innovation. Specifically, we are interested in how innovation inputs (e.g., R\&D expenditures) and outputs (e.g., citation-weighted value of patents and stock market value of patents) depend on the extent to which a firm is controlled by shareholders that have significant stakes in related firms and on the extent to which the innovation spills over to neighboring firms in the product market and technology space. As in our theoretical framework we study the economy-wide implications of common ownership and do not restrict ourselves to the study of any particular single industry. Unless otherwise stated, all of the data used for our estimations are from 1985 to 2015. Table~\ref{tab:sumstats} provides an overview of our summary statistics for the key variables. 

\subsection{Measures of Innovation}

To proxy for a firm's innovation $x_{i}$ in our theoretical model, we construct empirical innovation measures, denoted by $\mathit{INNOVATION}_{it}$, based on firm-level patent grants and citations from the database built by \cite{kogan2012technological}. This database has additions and corrections to the NBER patent data built by \cite{hall2001nber} from the official records of the United States Patent and Trademark Office (USPTO).\footnote{The database is available on Noah Stoffman’s website (\url{http://kelley.iu.edu/nstoffma}). More details on how to match patents and citations to the CRSP database can be found in the online appendix of \cite{kogan2012technological}.}

\begin{table}[!htbp]
\centering
\medskip
\resizebox{17.5cm}{!}{
\begin{tabular}{lrrrrrr}
\toprule
\textbf{Variables} & \multicolumn{1}{c}{\textbf{Obs}} & \textbf{Mean} & \textbf{Std. Dev.} & \textbf{10\%} & \textbf{50\%} & \textbf{90\%} \\
\midrule
\textit{\textbf{Innovation Variables}} &       &       &       &       &       &  \\
R\&D Expenditures &     28,596    &              245.14    &              953.00    &                 0.36    &     14.06    &        359.38    \\
ln(1+R\&D/Sales) &     28,596    &                  0.11    &                  0.20    &                  0.00    &        0.05    &             0.22    \\
TCW (Citation-weighted patents) &     28,631    &                87.45    &              418.96    &                  0.00    &        4.01    &        131.43    \\
TSM (Total stock market value of patents) &     28,631    &              930.73    &          5,867.81    &                  0.00    &        3.31    &        933.43    \\
          &       &       &       &       &       &  \\
\textit{\textbf{Proximity Measures \citep{bloom2013}}} &       &       &       &       &       &  \\
$\ln$(SPILLTECH) &     28,631    &                15.42    &                  1.01    &                14.12    &     15.53    &           16.60    \\
$\ln$(SPILLSIC) &     28,631    &                  4.49    &                  1.78    &                  2.32    &        4.67    &             6.65    \\
          &       &       &       &       &       &  \\
\textit{\textbf{Firm Characteristics}} &       &       &       &       &       &  \\
K/L (Capital labor ratio) &     28,200    &                55.22    &                77.16    &                11.05    &     31.45    &        116.95    \\
$\ln$(K/L) &     28,200    &                  3.52    &                  0.94    &                  2.41    &        3.44    &             4.75    \\
Sales &     28,631    &                4,080    &              14,265    &                      18    &         288    &           8,652    \\
Institutional Ownership &     27,761    &                0.463    &                0.295    &                0.080    &     0.464    &           0.833    \\
          &       &       &       &       &       &  \\
\textit{\textbf{Common Ownership}} &       &       &       &       &       &  \\
Kappa &     27,755    &                0.208    &                0.167    &                0.050    &     0.161    &           0.426    \\
\bottomrule
\end{tabular}
}
\protect\caption{Summary Statistics for Key Variables}
\label{tab:sumstats}
\end{table}


To measure innovation inputs, we use $\ln(1+R_{it}/\mathit{SALES}_{it})$ where $R_{it}$ is the level of inflation-adjusted R\&D expenditures and $\mathit{SALES}_{it}$ are the total sales of firm $i$ in year $t$ as reported in Compustat. Since many firms report zero values for R\&D expenditures in most years we follow \cite{branstetter2006stronger} by adding 1 to all observations of R\&D scaled by sales and taking the natural logarithm of the resulting values. 

To measure innovation outputs, we rely on two different measures that capture the scientific and economic value of innovation respectively. First, we use the number of citation-weighted patents $\mathit{TCW}_{it}$ given by
\begin{align}
\mathit{TCW}_{it} = \sum_{j \in P_{it}} \left(1 + \frac{C_{j}}{\bar{C}_{j}} \right)
\end{align}
where $P_{it}$ denotes the set of patents issued to firm $i$ in year $t$, $C_{j}$ is the number of forward citations to patent $j$, and $\bar{C}_{j}$ is the mean number of citations to patents granted in the same year as patent $j$. The innovation literature has argued that forward patent citations are a good indicator of the quality of the innovation and its scientific value \citep{hall2001nber}.

Second, we measure the private economic value of innovation \citep{hall2005,kogan2012technological} as proxied by stock market reactions following a patent issuance. Specifically, we use the measure of \cite{kogan2012technological} which estimates a firm $i$'s stock market reaction $\xi_{j}$ during the three-day announcement window following the issuance of the firm's patent $j$. \cite{kogan2012technological}  then sum up all the estimated values $\xi_{j}$ of patents $j$ that were granted to firm $i$ in year $t$ to construct the total stock market value of innovation $\mathit{TSM}_{it}$ generated by firm $i$ in year $t$: 
\begin{align}
\mathit{TSM}_{it} = \sum_{j \in P_{it}} \xi_{j}.
\end{align}

These two innovation outputs (forward patent citations and stock market value of patents) likely measure different aspects of quality. Whereas patent citations are more reflective of the scientific value of the innovation, the total stock market value measures the private economic value that is fully appropriated by the firm. For example, a patent may constitute only a minor scientific progress (and therefore generate few patent citations), but it may be particularly successful at limiting competition thereby generating significant profits for the issuing firm. 

% We follow common convention and use the natural logarithm of 1 plus the citation-weighted value of patents given by $\ln(1+\mathit{TCW}_{it})$. ACTUALLY WE DONT. Analogously to before, we use the natural logarithm of 1 + total stock market value of patents given by $\ln(1+\mathit{TSM}_{it})$.


\subsection{Measures of Technological and Product Market Proximity}

For our analysis we require two distinct measures of technological proximity and product market proximity between firms. We follow \citet{bloom2013} and \cite{lucking2018spillovers} by using $\mathit{SPILLTECH}$ and $\mathit{SPILLSIC}$ as measures of technological and product market spillovers. $\mathit{SPILLTECH}$ empirically measures the degree of technological spillovers $\sum_{j \neq i} \beta_{ij}$ and $\mathit{SPILLSIC}$ measures the degree of product market spillovers $\sum_{j \neq i} a_{ij}$ in our model. We briefly explain the specific construction of these measures below. For a thorough discussion including microeconomic foundations see \cite{bloom2013}. 

Following \cite{bloom2013} we construct $\mathit{SPILLTECH}$ and $\mathit{SPILLSIC}$ measures using both the Jaffe and the Mahalanobis distance.

\subsubsection{Jaffe Distance}

Denote the vector of the share of patents of firm $i$ in any given technology class by $T_i$. We construct the pool of technology spillover R\&D for firm $i$ in year $t$, $\mathit{SPILLTECH}_{it}$, as
\begin{equation}
\mathit{SPILLTECH}_{it} = \sum_{j \neq i} \mathit{TECH}_{ij} G_{jt}
\end{equation}
where $G_{jt}$ is the stock of R\&D and $\mathit{TECH}_{ij}$ is the uncentered correlation between all firm $i,j$ pairings and closely corresponds to the $\beta_{ij}$ parameter in our model. Following \cite{jaffe1988}, this measure is defined as
\begin{equation}
\mathit{TECH}_{ij} = \frac{T_{i} T^{'}_{j}}{(T_{i} T^{'}_{i})^{1/2} (T_{j} T^{'}_{j})^{1/2}}.
\end{equation}
We follow \cite{bloom2013} and \cite{lucking2018spillovers} by using all available years of data for the construction of the $\mathit{TECH}_{ij}$ measure.

Analogously, let $S_i$ denote the vector of the sales share of firm $i$ in any given four-digit industry. We construct the pool of product market spillover R\&D for firm $i$ in year $t$, $\mathit{SPILLSIC}_{it}$, as
\begin{equation}
\mathit{SPILLSIC}_{it} = \sum_{j \neq i} \mathit{SIC}_{ij} G_{jt}.
\end{equation}
This consists of the sum of the $\mathit{SIC}_{ij}$ uncentered correlations between all firm pairings $i$ and $j$ in an exactly analogous way to the technology closeness measure. This $\mathit{SIC}_{ij}$ measure closely corresponds to the $a_{ij}$ parameter in our model and is given by
\begin{equation}
\mathit{SIC}_{ij} = \frac{S_{i} S^{'}_{j}}{(S_{i} S^{'}_{i})^{1/2} (S_{j} S^{'}_{j})^{1/2}}.
\end{equation}
Following \cite{lucking2018spillovers} rather than pool across all years to construct a firm's industry sales share, we pool the previous five years of data. Pooling the industry segment data across all 30 years of our sample is problematic in this setting. Future industry sales shares are clearly endogenous as firm innovation and R\&D affect subsequent product market success. Past sales shares do not suffer from endogeneity but will be mismeasured if firms move in product space over time. We therefore use the five previous years of firm sales in order to (a) minimize reverse causality between firm outcomes and product market competition and (b) accurately measure the firm’s location in product market space at time $t$.


\subsubsection{Mahalanobis Distance}

Following \cite{bloom2013} and \cite{lucking2018spillovers} we also construct alternative versions of $\mathit{SPILLTECH}$ and $\mathit{SPILLSIC}$ using
the Mahalanobis distance metric which we denote by $\mathit{SPILLTECH}^{M}$ and $\mathit{SPILLSIC}^{M}$. These measures allow for spillovers between different technology classes and between different industries. In contrast, such spillovers across technology classes and four-digit industries are ruled out by the Jaffe metric which assumes full spillovers within the same class or industry and no spillovers otherwise. Complete detail on the definition and construction of the Mahalanobis measures is included in the online appendices of \cite{bloom2013} and \cite{lucking2018spillovers}. 

The Mahalanobis $\mathit{SPILLTECH}^{M}$ measure quantifies spillovers across technology class by using revealed preference. If two technologies are often located together in the same firm (for example, `computer input/output' and `computer processing'), then the distance between the technologies is smaller, so spillovers will be greater. The share of times the two technology classes are patented within the same firm proxies for this distance. The Mahalanobis $\mathit{SPILLSIC}^{M}$ measure is defined analogously for product market spillovers. When products within two industries are frequently sold within the same firm, the Mahalanobis $\mathit{SPILLSIC}^{M}$ measure infers that the distance between these two industries is small.

% Let $T = [T^{'}_{1},T^{'}_{2},...,T^{'}_{n}]$ denote the $(426 \times n)$ matrix where each column contains a firm $i$’s patent shares in the 426 technological classes for all $n$ firms. Define a normalized $(426,n)$ matrix $\Tilde{T} = [T^{'}_{1}/(T_{1} T^{'}_{1})^{1/2},...,T^{'}_{n}/(T_{n} T^{'}_{n})^{1/2}]$ in which each column is simply normalized by the firm’s patent share dot product. 
% % Define the $(n \times n)$ matrix $TECH = \Tilde{T}^{'}\Tilde{T}$. This matrix $TECH$ is just the standard \cite{jaffe1988} uncentered correlation measure between firms $i$ and $j$, in which each element is the measure T ECHij , exactly as defined in equation 1 above. 
% The Mahalanobis normed technology closeness measure is defined as T ECHM =
% Te0ΩTe. This measure weights the overlap in patent shares between firms by how close their different patents shares are to each other. The same patent class in different firms is given a weight of 1, and different patent classes in different firms are given a weight between 0 and 1, depending on how frequently they overlap within firms across the whole sample. Note that if Ω = I, then T ECHM = T ECH. Thus, if no patent class overlaps with any other patent class within the same firm, then the standard Jaffe (1986) measure is identical to the Mahalanobis norm measure. On the other hand, if some patent classes tend to overlap frequently within firms - suggesting they have some kind of technological spillover - then the overlap between
% firms sharing these patent classes will be higher.

\subsection{Measures of Common Ownership} 

To construct the ownership variables, we use Thomson Reuters 13Fs, which are taken from SEC regulatory filings by institutions with at least \$100m total assets under management. We augment this data by scraping SEC 13F filings following \cite{ben2016granular}, which resolves the issues of stale and omitted institutional reports, excluded securities, and missing holdings from 2000 onward.\footnote{This correction provides a full set of filings by institutional holders for the large universe of Compustat firms. For a detailed description of the problems with the original 13F database and the solutions provided see \url{https://wrds-www.wharton.upenn.edu/documents/752/Research_Note_-Thomson_S34_Data_Issues_mldAsdi.pdf}.} We describe the precise construction of the common ownership variables from these data in the following section. 

A limitation implied by this data source is that we do not observe the holdings of individual owners unless they are employed as officers of the company or serve on its board, in which case we complement these data with Execucomp. We assume that the remaining individual stakes of outsiders are relatively small and that in most cases they do not directly exert a significant influence on firm management. %Inspection of proxy statements of all firms in particular industries suggests that the stakes that individual outside shareholders own in large, publicly traded firms are rarely significant enough to substantially alter the measures of common ownership, even in the more prominent cases. For example, even Bill Gates's ownership of about 1.3\% of Microsoft's stock is small compared to the top five diversified institutional owners' holdings, which amount to more than 23\%.\footnote{As of 2020, Bill Gates no longer holds an officer or board position at Microsoft and is thus an outsider.} As a result, including or discarding the information on Bill Gates's holdings does not have a large effect on the measure of common ownership used.
The arising inaccuracies introduce measurement error and an attenuation bias toward zero in our regressions.


% The key explanatory variables are defined as follows. First, %MHHID
% $CO$ denotes two alternative measures of common ownership. For each company and year, we measure which fraction of the set of ``other'' companies is beneficially owned by the top five shareholders where the set of ``other'' companies is defined in various ways detailed below. We hereby aggregate the holdings of the top five owners for each firm, whereas we value-weight the other holdings by their market cap. Next, we rank-transform the common ownership variable.

% The alternative measure of common ownership is borrowed from Harford, Jenter, and Li (2011). In the $HJL$ measure, for each pair of companies and for each shareholder we first compute the following ownership ratio. \begin{equation*}
% \frac{\text{ownership of the shareholder in firm 1} \times \text{ownership of the shareholder in firm 2}}{\text{ownership of the shareholder in firm 1 + ownership of the shareholder in firm 2}}
% \end{equation*}
% We then sum this variable across all shareholders for each pair firm-pair. Finally, we aggregate the measure for each firm across all pairs, equal-weighting each pair.

To identify how common ownership influences the relationship between product market competition, technology spillovers, and innovation, we require a measure of common ownership. The existing literature provides several candidate measures of common ownership, the first of which is closely linked to the theoretical literature on common ownership, including our own model. 

From equation (\ref{eq:objective_function}), recall that the objective function of firm $i$ is given by
\begin{equation*}
\phi_{i} = \pi_{i} + \sum_{j \neq i} \kappa_{ij} \pi_j
\end{equation*}
where $\kappa_{ij}$ is the weight that firm $i$ places on firm $j$'s profits, $\pi_j$. The weighted sum of these profit weights $\kappa_{ij}$ across all the potential product market competitors of firm $i$ is the principal object of interest in the common ownership hypothesis \citep{backus2021common}. Our main measure of common ownership is $\kappa_{ij}$ between any firm pair $i$ and $j$ across the entire economy. We refer to the equal- or value-weighted average of the weights that the owners of firm $i$ place in year $t$ on the profits of the $n-1$ other firms in the economy as $\overline{\kappa}_{it}$ or simply ``kappa.'' More formally,
\begin{equation}
\overline{\kappa}_{it} = \frac{1}{n-1}\sum_{j \neq i} \kappa_{ij,t} \quad \textrm{or} \quad \overline{\kappa}_{it} = \frac{1}{\sum_{j \neq i} \omega_{jt}} \sum_{j \neq i} \kappa_{ij,t} \omega_{jt}
\end{equation}
where the weighting $\omega_{jt}$ is the stock market value of firm $j$ in year $t$. As in our theoretical model we exclude from our empirical analysis the small fraction of observations where $\overline{\kappa}_{it}$ exceeds 1 because these observations are indicative of incorrect or missing ownership data \citep{backus2021common}. Our results are essentially unchanged if we include these observations.

% Although the average profit weight $\overline{\kappa}_{i}$ is the leading measure for measuring common ownership and directly maps to the profit weights used in our theoretical analysis, it is important to verify that our empirical results are robust to using alternative measures of the extent to which a firm's most powerful shareholders care about competitor profits. 

% \cite{backus2021common} show that under proportional control (``one share, one vote''), each profit weight $\kappa_{ij}$ can further be decomposed into
% \begin{equation}
% \kappa_{ij} = \underbrace{\cos(\nu_i,\nu_j)}_\textrm{overlapping ownership} \cdot \underbrace{\sqrt{\frac{IHHI_j}{IHHI_i}}}_\textrm{relative IHHI}. 
% \end{equation}
% The first term is the cosine of the angle between the vector $\nu_i$ of ownership positions $\nu_{io}$ that owners (indexed by $o$) hold in firm $i$ and the corresponding vector $\nu_j$ for firm $j$. The second term is the ratio of the ``investor Herfindahl–Hirschman indices'' $IHHI_{i} = \sum_o \nu^{2}_{io}$ and $IHHI_{j} = \sum_o \nu^{2}_{jo}$ for the owners of firm $i$ and $j$.

% The cosine similarity captures the overlap in ownership and is the origin of the incentive to internalize the profits of another firm. Abstracting from the possibility of large short positions, ownership shares in $(i,j)$ are non-negative and therefore this similarity metric $\cos(\nu_i, \nu_j)$ is restricted to the $[0,1]$ interval. A cosine similarity of zero corresponds with no common owners while a cosine similarity of 1 corresponds to identical shareholding structure. Since this is an $L_2$ similarity measure, the metric puts more weight on large owners than small owners.\footnote{Note that the Hoberg-Phillips industry definitions also use a cosine similarity measure, which relies on the pairwise vectors of firm 10K product descriptions.} The second source of variation in common ownership profit weights comes from the ratio of the IHHI indices and ties the theory of common ownership to the notion that investor concentration drives a wedge between control rights and cash flow rights. All other things being equal, firms with concentrated investors will place more weight on their own profits and less weight on competitor profits.

% Ownership similarity is the symmetric component of the profit weight; if it increases, it will increase the objective functions of both firms in the industry. On the other hand, the relative shareholder concentration term is inherently asymmetric. To the extent that the asymmetric incentives of the profit-weight model are limited by legal restrictions or managerial behavior, empirically we may see the first-order effects of common ownership propagate through cosine similarity, as suggested by \cite{boller2019testing}. We therefore also use the equal-weighted average of the cosine similarity across all the $n-1$ competitors (indexed by $j$) of firm $i$ as a firm-specific measure for common ownership, which is given by
% \begin{equation}
% \overline{\cos}_{i} = \frac{1}{n-1}\sum_{j \neq i} \cos(\nu_i,\nu_j) \quad \textrm{or} \quad \overline{\cos}_{i} = \frac{1}{\sum_{j \neq i} \omega_{j}} \sum_{j \neq i} \cos(\nu_i,\nu_j) \omega_{j}.
% \end{equation}

% An alternative measure we employ is the average fraction of competitor shares held by the firm's top 5 shareholders, which we call the Top 5 shareholder measure. This is a model-free measure. In particular, this firm-specific measure for firm $i$ is defined as
% \begin{equation}
% \overline{Top5}_{i} = \frac{1}{n-1} \sum_{o}^{5} \sum_{i \neq j} \nu_{jo}  \quad \textrm{or} \quad \overline{Top5}_{i} = \frac{1}{\sum_{j \neq i} \omega_{j}} \sum_{o}^{5} \sum_{i \neq j} \nu_{jo} \omega_{j}
% \end{equation}
% where $\nu_{jo}$ is again the ownership share of firm $j$ accruing to shareholder $o$ who is one of the 5 largest owners of firm $i$, and $j$ indexes all of firm $i$'s competitors (of which there are $n-1$ for a given industry).


% Another established and popular measure of connectivity and ownership overlap between firms comes from \cite{anton2014connected}. It constructs a measure of common ownership that is the total value of stock held by all the common shareholders $o$ of two industry competitors $i$ and $j$, scaled by the total market capitalization of the two stocks $i$ and $j$. Specifically, this pair-level measure is
% \begin{equation*}
% AP_{ij} = \frac{\sum_{o} (S^{o}_{i}P_{i} + S^{o}_{j}P_{j})}{S_{i}P_{i}+S_{j}P_{j}}
% \end{equation*}
% where $S^{o}_{i}$ is the number of shares held by owner $o$ of firm $i$ trading at price $P_{i}$ with a total of  $S_{i}$ shares outstanding, and similarly for the stock of firm $j$. We use the weighted average across all $n-1$ industry competitors of firm $i$ and refer to it as the Anton-Polk (or AP) measure of common ownership:
% \begin{equation}
% \overline{AP}_{i} = \frac{1}{n-1}\sum_{j \neq i} AP_{ij} \quad \textrm{or} \quad \overline{AP}_{i} = \frac{1}{\sum_{j \neq i} \omega_{j}} \sum_{j \neq i} AP_{ij} \omega_{j}.
% \end{equation}

% We also use the modified cross-holdings measure from \cite{harford2011institutional} (henceforth the HJL measure), which accounts for the incentives of common investors during the merger of two firms. In their setting, the shareholders of a bidding firm are more likely to internalize the effect of paying a lower takeover premium on the target firm if they also own shares of the target. To capture this externality of common ownership, they estimate each investor's relative ownership stake in the target to that of the acquirer and aggregate these relative weights across investors in the bidding firm. Specifically, this pair-level measure is given by $HJL_{ij} = \sum_{o} \frac{\nu_{jo}}{\nu_{io}+\nu_{jo}}$. We use the (weighted) average of this measure across all industry competitors of firm $i$, given by
% \begin{equation}
% \overline{HJL}_{i} = \frac{1}{n-1}\sum_{j \neq i} HJL_{ij} \quad \textrm{or} \quad \overline{HJL}_{i} = \frac{1}{\sum_{j \neq i} \omega_{j}} \sum_{j \neq i} HJL_{ij} \omega_{j}.
% \end{equation}


\subsection{Other Variables}

Throughout our analysis we also use an additional set of control variables. First, $\ln(\mathit{SALES}_{it})$ is the natural logarithm of sales of the company where we adjust for inflation as in \citet{brav2014shareholder}. Second, $\ln(K_{it}/L_{it})$ is the capital-labor ratio, computed as the natural logarithm of the ratio of plant property equipment $K_{it}$ and the number of employees $L_{it}$ as in \cite{aghion2013innovation}, \cite{hall2001nber}, and \cite{gompers2001institutional}. Finally, we control for a firm's share of all of its institutional ownership as in \citet{aghion2013innovation} as this could also influence corporate innovation independent of the overlapping shareholdings of institutional investors.
%Third, $\ln(1+AGE_{it})$ is the natural logarithm of 1 +  firm $i$'s age as in \citet{brav2014shareholder}. Finally, $HHI_{zt}$ is the rank-transformed Herfindahl-Hirschman Index of market concentration defined as the sum of the squares of the market shares in industry $z$. We compute it using sales ($SALES_{it}$) defined at the industry level. Industries are defined at the 4-digit level using SIC codes from the CRSP data base, SIC codes from Compustat, or NAICS codes, alternatively. Note that we only require industry definitions to control for industry-time fixed effects, but we do not rely on them for our definition of product markets.


\section{Empirical Analysis}
\label{sec:results}

% we should probably control for the aghion institutional ownership variables!

We empirically study how corporate innovation depends on the degree to which the firms are commonly owned and how that relationship is affected by the spillovers on other firms in the technology and product market space. The theoretical model presented in Section~\ref{sec:model} illustrates that common ownership can have a positive or a negative effect on innovation, depending on parameters. Specifically, the model predicts that the correlation between common ownership and innovation increases with the level of technological spillovers but decreases the closer the firms are in product space.

% \subsection{Empirical Hypotheses}

% Whereas the model also makes predictions about the magnitude of the common ownership coefficient as a function of parameters, these predictions do not have a clear correspondence with the empirical methods we employ. We therefore do not offer tests that measure correlations between common ownership and innovation in a broad cross-section or time-series of the data, but focus on fixed-effect panel regressions instead. The results we obtain should therefore not be interpreted as globally valid relations between common ownership and innovation. Instead, we aim to provide locally valid estimates that help understand to which extent small variations in common ownership change innovation inputs and outputs, and how the intuition gained from the model is useful for understanding patterns of corporate innovation and common ownership in the data.

% I think that's a bunch of bullshit -- we should totally show the broad correlations as well!

% To be clear about what we are testing, however, we state the empirical null hypotheses. The two hypotheses we are testing correspond to the two corollaries presented in Section~\ref{sec:model}:

% Hypothesis 1:
% \begin{itemize}
% \item \emph{H0: Common ownership has no effect on innovation activity.}
% \item \emph{H1: Common ownership has a non-zero effect on innovation activity.}
% \end{itemize}

% Hypothesis 2:
% \begin{itemize}
% \item \emph{H0: The effect of Common Ownership on innovation does not depend on product market homogeneity or on technology proximity.}
% \item \emph{H1-spilltec: The effect of Common Ownership on innovation increases with the level of technological spillovers.}
% \item \emph{H1-spillsic: The effect of Common Ownership on innovation decreases with the level of product market spillovers.}
% \end{itemize}

\subsection{Empirical Methodology}

In our empirical analysis, we estimate for each of the three outcome variables (scaled R\&D, citation-weighted patents, stock market value of patents) how innovation depends on common ownership as well as the interactions of common ownership with product market spillovers and technology spillovers, controlling for known or suspected co-determinants of innovation such the size of the firm, capital intensity, and institutional ownership \citep{aghion2013innovation}. Our baseline regression is
\begin{multline}\label{eq:baseline}
\mathit{INNOVATION}_{it}= \alpha_{1} \cdot \mathit{CO}_{it} + \alpha_{2} \cdot \mathit{SPILLSIC}_{it} + \alpha_{3}  \cdot \mathit{SPILLTECH}_{it} \\ 
+ \alpha_{4} \cdot \mathit{CO}_{it} \cdot \mathit{SPILLSIC}_{it} + \alpha_{5} \cdot \mathit{CO}_{it} \cdot \mathit{SPILLTECH}_{it} \\
+ \alpha_{6} \cdot X_{it} + \sum_x \xi_{x} \cdot \eta_x + \varepsilon_{ijt}
\end{multline}
where firms are indexed by $i$, and years by $t$. $X_{it}$ is the vector of control variables $\ln(\mathit{SALES}_{it})$, $\ln(K_{it}/L_{it})$, and institutional ownership. $\eta_x$ with $x \in \{i,t\}$ are firm $i$, and year $t$ fixed effects. $CO_{it} = \overline{\kappa}_{it}$ measures to what extent the largest and most powerful shareholders of firm $i$ are also beneficial owners of other firms that are connected to firm $i$. Standard errors are clustered at the firm level.

We estimate OLS regressions for scaled R\&D expenditures and the stock market value of patents and negative binominal count data models for citation-weighted patents. The negative binomial regressions include a firm fixed effect that controls for the firm's average citation-weighted patents in the pre-sample period, as in \cite{blundell1999market}, \cite{bloom2013}, and \cite{lucking2018spillovers}, where the pre-sample period is defined as the five years before the firm enters the regression sample. 

Recall that firms influence each other because they benefit from any innovation activities of firm $i$ (technology spillovers) and/or because they are natural product market competitors of firm $i$ (product market spillovers). The principal coefficients of interest are therefore $\alpha_{4}$ and $\alpha_{5}$ which measure how the relationship between common ownership and innovation varies with product market and technology spillovers.


\subsection{Empirical Results}

We begin our analysis by examining the impact of common ownership and technology spillovers on innovation inputs (R\&D expenses) as the outcome variable. Table~\ref{tab:baselineRD} reports the results for estimation of equation (\ref{eq:baseline}) with the logarithm of one plus R\&D to sales ratio as the dependent variable. Across the different specifications we include firm and year fixed effects to control for otherwise potentially omitted time trends that would be correlated with changes in common ownership, as well as omitted controls that could be correlated with both firm-level differences in R\&D and common ownership and thus bias the common-ownership coefficient. Column 1 starts by replicating the results of \cite{lucking2018spillovers} as our baseline specification using the Jaffe proximity measures. Column 2 adds common ownership and shows that it is negatively correlated with innovation input. In columns 3 and 4, we include interaction terms between common ownership and our two proximity measures, and columns 5 and 6, which use the Mahalanobis proximity measures. Common ownership is associated with lower innovation input. We reject the null hypothesis that innovation is unrelated to a firm's common ownership. Columns 4 through 6 include controls for institutional ownership, as in \cite{aghion2013innovation}. We find a positive coefficient as they do, but it is statistically insignificant in our specifications.

\begin{table}[!htb]
\centering
\medskip
\resizebox{17.5cm}{!}{
\begin{tabular}{lcccccc}
\toprule
R\&D expenditure  & (1)   & (2)   & (3)   & (4)   & (5)   & (6) \\
$\ln(1+R_{it}/S_{it})$  & Jaf. & Jaf. & Jaf. & Jaf. & Mah. & Mah. \\
\midrule
$\mathit{CO}$    &       & -0.00151** & -0.0347** & -0.0347** & -0.00124* & -0.0697*** \\
      &       & (0.000736) & (0.0162) & (0.0164) & (0.000726) & (0.0215) \\
$\mathit{CO} \times \ln(\mathit{SPILLTECH})$ &       &       & 0.00247** & 0.00246** &       & 0.00486*** \\
      &       &       & (0.00108) & (0.00109) &       & (0.00143) \\
$\mathit{CO} \times \ln(\mathit{SPILLSIC})$ &       &       & -0.00111** & -0.00103** &       & -0.00224** \\
      &       &       & (0.000507) & (0.000513) &       & (0.000878) \\
$\ln(\mathit{SPILLTECH})$ & -0.0204*** & -0.0216*** & -0.0223*** & -0.0215*** & -0.0164* & -0.0176** \\
      & (0.00642) & (0.00651) & (0.00654) & (0.00668) & (0.00838) & (0.00841) \\
$\ln(\mathit{SPILLSIC})$ & 0.00468*** & 0.00442*** & 0.00470*** & 0.00489*** & 0.00215 & 0.00266 \\
      & (0.00134) & (0.00135) & (0.00135) & (0.00137) & (0.00219) & (0.00220) \\
\textit{Institutional Ownership} &       &       &       & 0.00426 & 0.00486 & 0.00471 \\
      &       &       &       & (0.00349) & (0.00350) & (0.00349) \\
      &       &       &       &       &       &  \\
Observations & 25,985 & 25,276 & 25,276 & 25,009 & 25,009 & 25,009 \\
Year FE & Yes   & Yes   & Yes   & Yes   & Yes   & Yes \\
Firm FE & Yes   & Yes   & Yes   & Yes   & Yes   & Yes \\
\bottomrule
\end{tabular}
}
\protect\caption{Sales-adjusted R\&D expenditure as a function of common ownership, technology spillovers, and product market spillovers
{\scriptsize{} \newline The table reports OLS coefficient estimates of equation (\ref{eq:baseline}) with the dependent variable $\ln(1+R_{it}/S_{it})$. Standard errors are clustered at the firm level. Variable definitions are described in Section~\ref{sec:data}. Industry sales in $t$ and in $t-1$ are also included as controls but are not shown in the table.}}
\label{tab:baselineRD}
\end{table}

Our primary interest, however, is with how the relation between common ownership and innovation varies with technology and product market spillovers. Columns 3, 4, and 6 include interaction terms between common ownership and our two measures of spillovers. Consistently throughout all the specifications, the estimated coefficient on the interaction term with technological spillovers $\mathit{SPILLTECH}$ is positive, whereas it is negative with product market spillover $\mathit{SPILLSIC}$. That is to say, in accordance with our theoretical analysis, the negative relation between common ownership and innovation inputs is more negative as the degree of product market spillovers increases. Conversely, the relationship between common ownership and innovation inputs becomes less negative and can even turn positive, when technology spillovers are large.

In aggregate, as can be seen in column 2, common ownership is negatively related to corporate innovation, but the estimated coefficient is quite small: increasing the average kappa from 0 (no common ownership) to 1 (equal profit weights on all other firms) is only associated with a -0.15\% (Jaffe) or -0.12\% (Mahalanobis) decrease in sales-adjusted R\&D expenditure. Thus, in the aggregate, the countervailing forces of technology and product market spillovers that pull the relationship between common ownership and corporate innovation in different directions essentially cancel each other out. However, this masks the significant heterogeneity in implied coefficient estimates of the relationship between common ownership and sales-adjusted R\&D expenditure. For about half of the firms (i.e., for those with relatively high technology and low product market spillovers) the relationship is positive whereas for the other half (low technology and high product market spillovers) it is negative.

\begin{table}[!tb]
\centering
\medskip
\resizebox{17.5cm}{!}{
\begin{tabular}{lcccccc}
\toprule
Citation-weighted patents & (1)   & (2)   & (3)   & (4)   & (5)   & (6) \\
$\mathit{TCW}_{it}$       & Jaf. & Jaf. & Jaf. & Jaf. & Mah. & Mah. \\
\midrule
$\mathit{CO}$    &       & 0.0476 & -6.085*** & -6.045*** & 0.125 & -6.520** \\
      &       & (0.138) & (2.162) & (2.189) & (0.147) & (2.675) \\
$\mathit{CO} \times \ln(\mathit{SPILLTECH})$ &       &       & 0.465*** & 0.466*** &       & 0.519*** \\
      &       &       & (0.154) & (0.156) &       & (0.185) \\
$\mathit{CO} \times \ln(\mathit{SPILLSIC})$ &       &       & -0.237*** & -0.233** &       & -0.346*** \\
      &       &       & (0.0919) & (0.0928) &       & (0.134) \\
$\ln(\mathit{SPILLTECH})$ & 0.133*** & 0.133*** & 0.0400 & 0.0397 & 0.174*** & 0.0676 \\
      & (0.0475) & (0.0475) & (0.0566) & (0.0567) & (0.0614) & (0.0707) \\
$\ln(\mathit{SPILLSIC})$ & -0.0130 & -0.0127 & 0.0367 & 0.0372 & -0.0237 & 0.0521 \\
      & (0.0257) & (0.0256) & (0.0317) & (0.0319) & (0.0371) & (0.0459) \\
\textit{Institutional Ownership} &       &       &       & 0.137 & 0.128 & 0.137 \\
      &       &       &       & (0.0952) & (0.0936) & (0.0953) \\
      &       &       &       &       &       &  \\
Observations & 24,683 & 24,683 & 24,683 & 24,487 & 24,487 & 24,487 \\
Year FE & Yes   & Yes   & Yes   & Yes   & Yes   & Yes \\
Firm FE & Yes   & Yes   & Yes   & Yes   & Yes   & Yes \\
\bottomrule
\end{tabular}
}
\protect\caption{Citation-weighted measure of patents as a function of common ownership, technology spillovers, and product market spillovers 
{\scriptsize{} \newline The table reports coefficient estimates of equation (\ref{eq:baseline}) with the dependent variable $\mathit{TCW}_{it}$ using a negative binominal count data model. Standard errors are clustered at the firm level. Variable definitions are described in Section~\ref{sec:data}. Industry sales in $t$ and in $t-1$ are also included as controls but are not shown in the table.}}
\label{tab:baselineTCW}
\end{table}

We now turn to the empirical relation between common ownership and innovation outputs. Table~\ref{tab:baselineTCW} reports the results for the citation-weighted value of patents held by a firm using a negative binomial count data model. On average, the citation-weighted value of patents is weakly positively correlated with common ownership as shown in column 2, but the estimated coefficient is statistically insignificant. 
% Even a very large increase in common ownership from 0 (no common ownership) to 1 (full common ownership) is only associated with a 4.9\% ($=e^{0.0476}-1$) increase in citation-weighted patents. 
In other words, on average across our entire sample, common ownership and corporate innovation are not particularly strongly related. However, this is because the interactions of common ownership with technology and product market spillovers cancel each other out in the aggregate, as the next columns show.

In particular, once we include the interaction terms between common ownership and the two spillover measures in columns 3, 4, and 6, the coefficient on common ownership is consistently negative. As before and in accordance with our theoretical predictions, we find that this negative relationship is larger in magnitude when product market spillovers $\mathit{SPILLSIC}$ are larger, but smaller in magnitude and even positive when technology spillovers $\mathit{SPILLTECH}$ are larger. Because there is considerable heterogeneity of these spillovers across industries and firms, this leads to vastly different effects of common ownership on corporate innovation for different firms. For example, for a firm at the 75th percentile of technology spillovers (16.13) and the 25th percentile of product market spillovers (3.53), an interquartile range increase of common ownership from the 25th percentile (0.092) to the 75th percentile (0.277) is associated with a 12.5\% increase in citation-weighted patents.\footnote{The exact calculation of this estimate comes from column 4 of Table~\ref{tab:baselineTCW} and is given by $\exp[(-6.045 + 0.466 \times 16.13 - 0.237 \times 3.53) \times (0.277 - 0.092)] - 1 = 0.125$.} In contrast, the same increase in common ownership implies a decrease of -8.4\% in citation-weighted patents for a firm at the 25th percentile of technology spillovers and 75th percentile of product market spillovers. 

% Another way is to just use the distribution of implied coefficient estimates
% For TCW:
% 10% -7.72%
% 25% -2.55%
% 75% +6.85%
% 90% +11.20%   
% 
% For TSM:
% 10% +7.03%
% 25% +10.66%
% 75% +17.52%
% 90% +20.51% 


\begin{figure}[!t]
\centering 
\includegraphics[width=17.5cm]{tables/fig:coefs_tcw.png}
\captionsetup{labelfont=bf}
\caption{\textbf{Heterogeneity of the Relationship between Common Ownership and Citation-weighted Patents $\mathit{TCW}$} \\This figure plots the distribution of how an increase in common ownership from the 25th to the 75th percentile changes citation-weighted patents $\mathit{TCW}$ taking into account firm-specific levels of technology and product market spillovers (averaged across years). Negative coefficient estimates are shown in \textcolor{red}{red} and positive ones in black.}
\label{fig:coefs_tcw}
\end{figure}

We illustrate the significant heterogeneity of the relationship between common ownership and corporate innovation in Figure~\ref{fig:coefs_tcw}. Specifically, we plot the aggregate coefficient estimate for common ownership taking into account how technology and product market spillovers vary across firms. For each firm, we compute its average value of $\mathit{SPILLTECH}$ and $\mathit{SPILLSIC}$ across all the years of our sample. We use these values in conjunction with the coefficient estimates of $\mathit{CO}$, $\mathit{CO} \times \ln(\mathit{SPILLTECH})$, and $\mathit{CO} \times \ln(\mathit{SPILLSIC})$ from column 4 of Table~\ref{tab:baselineTCW} to plot how a firm-level interquartile range increase of common ownership from the 25th to the 75th percentile changes the percentage of citation-weighted patents. 


\begin{table}[!tb]
\centering
\medskip
\resizebox{17.5cm}{!}{
\begin{tabular}{lcccccc}
\toprule
Patent stock market value & (1)  & (2)   & (3)   & (4)   & (5)   & (6) \\
$\ln(1+\mathit{TSM}_{it})$       & Jaf. & Jaf. & Jaf. & Jaf. & Mah. & Mah. \\
\midrule
$\mathit{CO}$    &       & 0.581*** & -2.779*** & -2.883*** & 0.692*** & -3.868*** \\
      &       & (0.0680) & (0.964) & (1.002) & (0.0797) & (1.318) \\
$\mathit{CO} \times \ln(\mathit{SPILLTECH})$ &       &       & 0.206*** & 0.219*** &       & 0.263*** \\
      &       &       & (0.0701) & (0.0729) &       & (0.0921) \\
$\mathit{CO} \times \ln(\mathit{SPILLSIC})$ &       &       & 0.0418 & 0.0446 &       & 0.0463 \\
      &       &       & (0.0523) & (0.0532) &       & (0.0809) \\
$\ln(\mathit{SPILLTECH})$ & 0.0785 & 0.0625 & -0.0122 & 0.0327 & -0.0530 & -0.136 \\
      & (0.0907) & (0.0902) & (0.0929) & (0.0936) & (0.117) & (0.120) \\
$\ln(\mathit{SPILLSIC})$ & 0.0508* & 0.0541** & 0.0443 & 0.0413 & 0.0775* & 0.0652 \\
      & (0.0268) & (0.0268) & (0.0272) & (0.0274) & (0.0421) & (0.0419) \\
\textit{Institutional Ownership} &       &       &       & 0.367** & 0.350** & 0.357** \\
      &       &       &       & (0.147) & (0.143) & (0.146) \\
      &       &       &       &       &       &  \\
Observations & 24,694 & 24,694 & 24,694 & 24,495 & 24,495 & 24,495 \\
Year FE & Yes   & Yes   & Yes   & Yes   & Yes   & Yes \\
Firm FE & Yes   & Yes   & Yes   & Yes   & Yes   & Yes \\
\bottomrule
\end{tabular}
}
\protect\caption{Stock market value of patents as a function of common ownership, technology spillovers, and product market spillovers 
{\scriptsize{} \newline The table reports OLS coefficient estimates of equation (\ref{eq:baseline}) with the dependent variable $\ln(1+\mathit{TSM}_{it})$. Standard errors are clustered at the firm level. Variable definitions are described in Section~\ref{sec:data}. Industry sales in $t$ and in $t-1$ are also included as controls but are not shown in the table.}}
\label{tab:baselineTSM}
\end{table}

Similar patterns emerge in Table~\ref{tab:baselineTSM} which reports the coefficient estimates for the relationship between the total stock market value of patents and common ownership. On average, the total stock market value of patents is now significantly positively correlated with common ownership as shown in columns 2 and 5. However, once we include the interaction terms between common ownership and the two spillover measures in columns 3, 4, and 6, the coefficient on common ownership is consistently negative. As predicted by our theoretical framework, we find that this negative relationship is weaker when technology spillovers $\mathit{SPILLTECH}$ are larger. That is to say, as before, depending on the importance of technological spillovers across the universe of firms, common ownership can either be negatively or positively related to corporate innovation. However, unlike for R\&D expenditure and citation-weighted patents, the coefficient estimate for the interaction of common owners and product market spillovers $\mathit{SPILLSIC}$ is not  statistically significant. The coefficient on institutional ownership is positive and significant, in line with the results of \cite{aghion2013innovation}. 


The relative importance of technology spillovers leads to significant heterogeneity in the relationship between common ownership and innovation as can be seen in Figure~\ref{fig:coefs_tsm}. However, in contrast to our other two innovation measures, the vast majority of implied coefficient estimates is positive and only a small number is negative. 


The finding of small effects on R\&D inputs but economically large effects (in some firms) on R\&D outputs is  consistent with the findings of \cite{li2020venture}, who find that common ownership by venture capitalists in the pharmaceutical industry reduces duplication of R\&D and thus increases innovation efficiency.

\begin{figure}[!ht]
\centering 
\includegraphics[width=17.5cm]{tables/fig:coefs_tsm.png}
\captionsetup{labelfont=bf}
\caption{\textbf{Heterogeneity of the Relationship between Common Ownership and the Market Value of Patents $\mathit{TSM}$} \\This figure plots the distribution of how an increase in common ownership from the 25th to the 75th percentile changes the market value of patents $\mathit{TSM}$ taking into account firm-specific levels of technology and product market spillovers (averaged across years). Negative coefficient estimates are shown in \textcolor{red}{red} and positive ones in black.}
\label{fig:coefs_tsm}
\end{figure}

Taken together, we find strong and fairly consistent support for the model's theoretical predictions. First, there exists an empirically ambiguous relationship between common ownership and innovation that can be either positive or negative on average. Second, the innovation-reducing effect of common ownership increases with the degree of product market spillovers. Third, technology spillovers increase the innovation-enhancing effect of common ownership. As predicted by the theoretical framework, the overall effect of common ownership on corporate innovation crucially depends on the relative strength of product market business stealing incentives and of technological spillovers between firms and differs markedly across firms. 

One limitation of the results presented so far is that, in the theoretical model, ownership is taken as an exogenously given parameter, whereas the data we rely on for our empirical estimation of common ownership effects of innovation are the result of a process involving choices (and hence selection) of ownership. This could potentially lead to an endogeneity problem that challenges a causal interpretation of the panel correlations we estimate. That said, we find it difficult to formulate, even informally, a simple economic model that would give rise to the two opposing effects we measure, but not allow for a causal interpretation. Take, for example, a perfectly ``passive'' investor holding the market portfolio as a benchmark. Some types of active investors might pursue a strategy of underweighting industry competitors while overweighting technologically related firms, and at the same time they might have a preference for more innovation (for reasons unrelated to their portfolio choice). Other types of active investors might deliberately overweight industry competitors and sell other holdings while pushing firms to reduce innovation---again for reasons unrelated to their portfolio choice. We are not aware of an economic rationale that links these portfolio choices with a preference for low or high innovation. Or perhaps firms with high innovation activity attract active shareholders that tend to underweight industry competitors in their investment strategy, whereas firms with low innovation activities attract active investors that specialize in holding industry competitors. Again, we are not aware of an economic rationale that could give rise to such a relationship. 

In contrast, the economic model proposed in the present paper provides a significantly simpler and economically intuitive explanation for the empirical patterns we observe. We are not aware of any alternative theory that would give rise to an endogeneity problem and suggest a different interpretation of the empirical evidence. In sum, any causal interpretation of the empirical correlations we documented thus far would come from combining theory with a set of nuanced patterns in the data, rather than from empirical observations alone.

% The story is more complicated for the comparative statics with respect to product market spillovers, however: whereas more commonly owned firms seem to spend more money on R\&D in markets with greater levels of product market spillovers, these expenditures do not seem to get translated into innovation outputs in an effective way. To the contrary, innovation outputs are less related to common ownership the higher the level of product market spillovers. Hence, given how inefficient the greater R\&D expenditures seem to be spent, there is unlikely to be a business stealing effect from these resources spent.

% Taken together, this evidence points towards a positive relationship between innovation and common ownership both for innovation inputs, and a negative relationship between innovation and common ownership for innovation outputs. The evidence is also consistent generally with corollary 2 of the theory: the effect of common ownership on innovation is increasing in SPILLTEC and decreasing in SPILLSIC.

% Table XX takes our analysis one step further by investigating how the above results differ when we use an exogenous measure of the spillover measures. Following \cite{bloom2013}, we now employ an instrumented version of the spillover measures. The reason is the following: in our baseline regressions, the dependent variable is R\&D, and the independent variables include SPILLTEC and SPILLSIC, which are computed using R\&D. We use as instrument the supply side shocks from tax-induced changes to the user cost of R\&D capital. Details for this identification approach can be found in \cite{bloom2013}.

% Perhaps common owners are ``lazy'' shareholders who monitor their portfolio firms insufficiently, and thus allow them to squander R\&D resources? We investigate this and further questions in ongoing research.
% examine this with governance measures?
% Table \ref{tab:interactions} reports results from a more formal analysis of how interactions of competition, technological spillovers and product market spillovers affect innovation. 

%\section{IV Strategy and Results}
%\label{sec:iv}

\section{Conclusion}
\label{sec:conclusion}

In this paper we show that common ownership can increase innovation when technological spillovers are sufficiently large relative to product market spillovers. On the other hand, common ownership can also decrease innovation because common owners would like to discourage business stealing between commonly owned companies that compete in product markets against each other. The ambiguity of the theoretical prediction thus poses an interesting empirical question about the sign and magnitude of the effect of common ownership on innovation. We use our theoretical model's predictions to investigate how the relationship between common ownership and innovation depends on the relative strength of technological and product market spillovers. Consistent with the model's theoretical predictions, we find that common ownership has a positive effect on innovation inputs and outputs whenever innovation spillovers to other firms are relatively large compared to the firms' distance in the product market space and a negative effect if the product market spillovers dominate.

% Our theoretical and empirical results are consistent with known relationships between ownership, firm size and innovation and their time-series trends. First, \cite{bernstein2015does} provides evidence that going public reduces innovation activity. Second, \cite{itenberg2015} finds that innovation has shifted from bigger to smaller firms over the past decades. Because initial public offerings increase common ownership and common ownership has significantly increased over the past decades, particularly for larger firms \citep{backus2021common}, our analysis also suggests an explanation for these trends. Whether common ownership is indeed the variable that causally links these trends is an interesting question for future research

Our findings inform an active debate on whether welfare-enhancing effects of common ownership outweigh the previously empirically documented negative effects of common ownership on firms' incentives to compete. Because a positive effect on innovation---which we model as an efficiency increase in this paper---is a necessary condition for common ownership to positively affect welfare \citep{lopez2016cross}, findings of positive innovation effects are a necessary ingredient in the argument against regulatory interventions on horizontal common ownership links that have competitive effects. The more nuanced insight, however, is that antitrust and innovation policy should distinguish between common ownership of horizontal competitors and common ownership of technologically and perhaps vertically related firms. Previous literature indicates that the former weakens competition and, as we show, also reduces innovation. Our theoretical analysis and empirical results suggest that the latter promotes innovation and potentially even increases total welfare.

%Taking our results at face value would suggest that policy proposals aiming to reduce anticompetitive horizontal common ownership in favor of greater common ownership of otherwise related firms, such as \cite{posner2016}, may have the effect of strengthening innovation.

%A more subtle insight, however, is that the cross-sectional differences across firms and industries in the effect of common ownership on innovation suggest that antitrust policy with respect to common ownership (as recently discussed by \cite{posner2016}) may benefit from targeting some industries more intensely than others. The presence of substantial technology spillovers mitigates and can even outweigh the anticompetitive effects of common ownership arising from business stealing incentives.

\setstretch{0.96}
\bibliographystyle{aer}
\bibliography{references}
\normalsize

%\newpage
\setstretch{1.5}
% \renewcommand{\thetable}{\arabic{table}}
% \input{figstables.tex}

\newpage
\appendix
\section{Proofs and Additional Theoretical Results}
\label{sec:appendix}
\input{appendix.tex}

%\newpage{}
%
%\section*{Appendix B: Additional Empirical Results}
%
%\renewcommand{\thetable}{\Roman{table}}
%\setcounter{table}{0}
%\setcounter{figure}{0}
%\renewcommand{\tablename}{Table A.}
%\renewcommand{\figurename}{Figure A.}
%
%\input{appendixtables}

\end{document}
